\chapter{黏性流体动力学基础}

\begin{yijubox}
E1 \texttt{8-1考情分析.pptx} slide33：第八章黏性流体动力学基础＝\textbf{选择填空题考点}。\\
E1 slide33 同时把第七、八章一并归入"概念为主"的收尾区；E2 \texttt{9-1考点分析.pptx}：第六至第八章只有少量填空/判断/选择，\textbf{掌握基本概念＋公式套用}即可。\\
E1 slide16：选择/判断考\emph{概念辨析}，填空考\emph{名词定义与小计算}。\\
E4 证据一：2015 年试题计算题含"\textbf{平板摩擦阻力比}"；E4 证据二：2015 年试题填空题
"层流断面平均速度＝0.5 倍最大速度、$\Rey$ 判流态"。\\
\textbf{本章结论}：本章\textbf{不出大题}，命题集中在三种题型——
① 名词定义型填空（边界层、边界层厚度、位移厚度、动量损失厚度、雷诺应力等）；
② 概念辨析型选择/判断（边界层内压强分布、分离的条件、$\Rey$ 判据的取法）；
③ 公式代入型小计算（$\delta$、$C_f$、$C_D$、阻力比）。
复习策略＝\textbf{把定义记死、把公式认准、把易错点背下来}。篇幅重心应放在 8.2（概念）、8.7（简答）、8.1（概念）三节。
\end{yijubox}

\begin{tishi}
\textbf{与第 5 章的分工（本讲义不重复展开）：}雷诺实验、两种流态与临界雷诺数、
圆管层流的速度分布（抛物线、$u_{\max}=2v$）、圆管层流沿程水头损失、湍流的沿程阻力系数，
均已在\textbf{第 5 章}讲过，本章只在必要的衔接处一句带过，一律注明"详见第 5 章"。
本章的新东西是：\emph{黏性流体的运动微分方程（N-S 方程）}、\emph{边界层理论}、\emph{湍流的时均方程与模型}。
\end{tishi}

\section{纳维-斯托克斯方程}

\begin{kaodian}{3}{黏性流体中的应力状态与本构关系}
\textbf{是什么：}理想流体运动时可以存在相对运动，但不具有抵抗剪切变形的能力，任意面上的应力只有法向力；
\textbf{黏性流体}具有抵抗剪切变形的能力，任意面上的应力\textbf{既有法向应力也有切应力}。

\textbf{应力记号（必考，注意下标的读法）：}两个下标表示清楚"作用面"与"投影方向"——
\textbf{第一个下标＝作用面的法线方向}，\textbf{第二个下标＝应力分量的投影方向}。
如 $x$ 面上的合应力 $\vec{\tau}_x=\tau_{xx}\vec{i}+\tau_{xy}\vec{j}+\tau_{xz}\vec{k}$。
三个坐标面上的 9 个分量组成\textbf{应力张量}；由剪力互等定理
$\tau_{xy}=\tau_{yx}$、$\tau_{xz}=\tau_{zx}$、$\tau_{yz}=\tau_{zy}$，应力矩阵是\textbf{对称矩阵}，
9 个分量中只有 \textbf{6 个独立}。

\textbf{平均压强：}\emph{理想流体}中无切应力且三个法向应力相等，$\tau_{xx}=\tau_{yy}=\tau_{zz}=-p$；
\emph{黏性流体}中过一点任意三个相互垂直面上的法向应力之和是不变量，定义
\[ p=-\frac{\tau_{xx}+\tau_{yy}+\tau_{zz}}{3} \]
为\textbf{平均压强}。注意：黏性流体中该和虽是不变量，但各法向应力本身\emph{一般不等于} $-p$。

\begin{gongshi}{广义牛顿内摩擦定律（本构关系，8.1）}
% 原为三条式子用 \quad 横向排在同一个 \[ \] 里，实测该盒超出右边距 5.18033pt；
% 改为 \begin{aligned} 纵向排列，每行只放一条式子，彻底消除溢出。
\[ \begin{aligned}
   \tau_{xy}&=\tau_{yx}=\mu\left(\frac{\partial u_x}{\partial y}+\frac{\partial u_y}{\partial x}\right),\\
   \tau_{yz}&=\tau_{zy}=\mu\left(\frac{\partial u_y}{\partial z}+\frac{\partial u_z}{\partial y}\right),\\
   \tau_{xz}&=\tau_{zx}=\mu\left(\frac{\partial u_x}{\partial z}+\frac{\partial u_z}{\partial x}\right)
   \end{aligned} \]
\[ \sigma_{xx}=-p+2\mu\frac{\partial u_x}{\partial x},\quad
   \sigma_{yy}=-p+2\mu\frac{\partial u_y}{\partial y},\quad
   \sigma_{zz}=-p+2\mu\frac{\partial u_z}{\partial z} \]
\textbf{含义：}把切应力与\emph{角变形速度}（两个方向偏导数的和）、法向应力与\emph{线变形率}线性联系起来，
是牛顿流体在三维情形下对牛顿内摩擦定律 $\tau=\mu\,du/dy$ 的推广。法向应力中 $-p$ 是各向同性的压强部分，
$2\mu\,\partial u_x/\partial x$ 是因线变形率不同而附加的黏性法向应力。\\
\textbf{教材式号：}（教材式 8.4）切应力普遍表达式；（教材式 8.3）牛顿内摩擦定律 $\tau=\mu\,\dd\theta/\dd t$；
（教材式 8.5）$p_{xx}=p+p'_{xx}$ 等（平均压强＋附加压应力）；（教材式 8.6）附加压应力；
（教材式 8.7）压应力与线变形速度的关系式（书 p154--155）。\\
\textbf{适用条件：}\textbf{牛顿流体}、连续介质、不可压缩（或采用斯托克斯假设：体积黏度 $\mu'=0$，
即只保留 $-p$ 而不加 $\mu'\,\mathrm{div}\,\bv$ 项）。
\end{gongshi}

\textbf{以应力表示的流体运动微分方程}（沿 $x$ 方向）：
\[ X+\frac{1}{\rho}\left(\frac{\partial\tau_{xx}}{\partial x}+\frac{\partial\tau_{yx}}{\partial y}+\frac{\partial\tau_{zx}}{\partial z}\right)=\frac{\dd u_x}{\dd t} \]
它由\emph{牛顿第二定律}直接写出，但\textbf{方程组不封闭}（未知的应力分量多于方程数），
必须补入本构关系才能求解——这正是下一步推导 N-S 方程的出发点。
\textbf{教材式号：}（教材式 8.1）为该方程的中间推导式；（教材式 8.2a）（教材式 8.2b）即"以应力表示的黏性流体运动微分方程"（书 p153）。

\textbf{常考题型：}填空（"应力分量 $\tau_{xy}$ 中第一个下标表示\underline{\ 作用面法线方向\ }"）；
选择（"黏性流体中任意一点的法向应力"）。
\textbf{重要度 \xstarfull{3}}\quad\yiju{E1 slide16 填空考"名词定义"；E2"掌握基本概念"；E6 教材 8.1 节}\quad\nandu{易}

\begin{banfa}
① \textbf{下标口诀："先面后向"}———先写面（法线方向），后写向（投影方向）。
② \textbf{数一数就得分：}$3\times3=9$ 个分量，对称性砍掉 3 个$\to$独立 6 个。
③ 判断题见到"黏性流体中三个法向应力都等于 $-p$"$\to$\textbf{错}（那是理想流体的结论）；
见到"三个法向应力之和是不变量"$\to$\textbf{对}。
\end{banfa}
\end{kaodian}

\begin{kaodian}{3}{纳维-斯托克斯方程的建立与各项物理意义}
\begin{gongshi}{不可压缩牛顿流体的 N-S 方程（8.1，分量形式）}
\[ \rho\frac{\dd u_x}{\dd t}=\rho X-\frac{\partial p}{\partial x}+\mu\left(\frac{\partial^2 u_x}{\partial x^2}+\frac{\partial^2 u_x}{\partial y^2}+\frac{\partial^2 u_x}{\partial z^2}\right) \]
（$y$、$z$ 方向同理）或除以 $\rho$ 写成单位质量形式：
\[ X-\frac{1}{\rho}\frac{\partial p}{\partial x}+\nu\left(\frac{\partial^2 v_x}{\partial x^2}+\frac{\partial^2 v_x}{\partial y^2}+\frac{\partial^2 v_x}{\partial z^2}\right)=\frac{\dd v_x}{\dd t} \]
\textbf{含义：}把\emph{广义牛顿内摩擦定律}代入\emph{以应力表示的运动微分方程}，
消去应力分量后得到的速度场控制方程。物理上就是"对流体微团应用牛顿第二定律"，
左边是惯性，右边是作用在微团上的全部外力。\\
\textbf{教材式号：}（教材式 8.8a）（教材式 8.8b）（教材式 8.8c）（书 p156）。
\textbf{教材原话（书 p156）：}"上式即为不可压缩均质黏性流体的运动微分方程，即纳维-斯托克斯方程，简称 N-S 方程。
如果流体是理想流体，上式则成为理想流体的运动微分方程；如果流体为静止流体，上式则成为欧拉平衡微分方程。
所以，N-S 方程是不可压缩均质流体的\textbf{普遍方程}。"
教材同时说明："N-S 方程是二阶非线性偏微分方程组……只有在某些简单的或特殊的情况下，才能求得精确解。"（书 p156）\\
\textbf{适用条件：}牛顿流体＋不可压缩＋均质（常黏度 $\mu$）。
\end{gongshi}

\begin{tishi}
\textbf{订正（原表述 $\to$ 新表述）：}本公式盒原写"写成上式的形式时\textbf{默认层流}，此时 N-S 方程是严格成立的方程"，
容易被读成"N-S 方程\textbf{只适用于层流}"。按教材订正为：N-S 方程是"不可压缩均质流体的\textbf{普遍方程}"（书 p156），
它与"理想流体的运动微分方程""欧拉平衡微分方程"是\emph{普遍与特殊}的关系；
对湍流同样成立——教材原话："理论上描述\textbf{瞬时流动}的 N-S 方程\textbf{同样适用于湍流}"（书 p176）。
准确说法是：\textbf{N-S 方程对层流与湍流的瞬时流动都成立}，只是"求解湍流瞬时流动的全部过程，既不必要也不可能"，
故改用\emph{时均化}方法得到雷诺方程（教材原话，书 p176；见 8.8 节）。
\end{tishi}

\begin{gongshi}{不可压缩 N-S 方程的矢量形式（必背）}
\[ \rho\left(\frac{\partial \bv}{\partial t}+(\bv\cdot\grad)\bv\right)=\rho\bF-\grad p+\mu\grad^2\bv \]
\textbf{含义：}等式左边是单位体积流体的\textbf{惯性项}（$\partial\bv/\partial t$ 为当地加速度、$(\bv\cdot\grad)\bv$ 为迁移加速度）；
右边依次是\textbf{质量力项}、\textbf{压力项}、\textbf{黏性项}。\\
\textbf{教材式号：}教材把 N-S 方程写成三个分量式并编为（教材式 8.8a）～（教材式 8.8c）（书 p156），
\textbf{矢量形式教材未单独编号}，本讲义补出矢量式以便记忆。\\
\textbf{适用条件：}不可压缩牛顿流体；$\mu=$常数；需要与连续性方程 $\grad\cdot\bv=0$ 联立封闭。
\end{gongshi}

\textbf{方程中四项的物理意义（选择题高频，务必逐项会说）：}
\begin{center}
\begin{tabularx}{\textwidth}{@{}lXX@{}}
\toprule
项 & 名称 & 物理意义（作用在单位体积流体上） \\
\midrule
$\rho\dfrac{\dd \bv}{\dd t}$ & 惯性项 & 使微团产生加速度所需的力；其中 $(\bv\cdot\grad)\bv$ 是\textbf{非线性}的来源 \\
$-\grad p$ & 压力项 & 由压强梯度产生的\emph{各向同性}的面力，方向由高压指向低压 \\
$\mu\grad^2\bv$ & 黏性项 & 分子动量扩散产生的\emph{耗散性}切力，把动能不可逆地转为热能 \\
$\rho\bF$ & 质量力项 & 外力场（重力、惯性力等）对单位体积流体施加的\textbf{非接触力} \\
\bottomrule
\end{tabularx}
\end{center}

\textbf{为什么难解（必背结论）：}N-S 方程是\textbf{二阶非线性偏微分方程}，
非线性来自惯性项中的迁移加速度 $(\bv\cdot\grad)\bv$；
再加上\textbf{边界条件难以用解析式表达}，除极少数简单情形外\textbf{得不到解析解}，
工程上只能依靠数值计算（CFD）与湍流模型。

\begin{jielun}
$\Rey$ 数可理解为惯性项与黏性项的量级之比 $\Rey\sim\rho U^2/l \div \mu U/l^2=\rho Ul/\mu$。
$\Rey$ 很小时黏性项占优（层流、可求解析解）；$\Rey$ 很大时惯性项占优，非线性不可忽略，
这就是高雷诺数流动必须另辟蹊径（边界层理论、湍流模型）的根本原因。
\end{jielun}

\textbf{常考题型：}选择（"N-S 方程中 $(\bv\cdot\grad)\bv$ 项的物理意义"）；判断（"N-S 方程只适用于层流"）；
填空（"N-S 方程不能求解析解的原因是\underline{\ 惯性项非线性\ }"）。
\textbf{重要度 \xstarfull{3}}\quad\yiju{E1 slide33 第八章＝"选择填空题考点"；E1 slide16 选择/判断考"概念辨析"；E2"掌握基本概念"}\quad\nandu{中}

\begin{banfa}
① \textbf{"四项对号入座"速记：}\emph{惯性—压力—黏性—质量力}，读式子时按 $\rho\bv$、$p$、$\mu$、$F$ 四个符号扫一遍就不会漏。
② 判断"N-S 方程只适用于层流"$\to$\textbf{错}。它本身是黏性流体的运动方程；
写成 $\rho\frac{\dd \bv}{\dd t}=\rho\bF-\grad p+\mu\grad^2\bv$ 时对\emph{瞬时量}成立，
湍流时把 $\bv$、$p$ 替换为瞬时量后再时均化即可（见 8.8 节）。
③ "为什么不能求解析解"三句话答题法：\textbf{①惯性项非线性 ②边界条件复杂 ③高 $\Rey$ 下解不稳定（湍流）}。
\end{banfa}

\begin{yicuo}
\textbf{（硬性纠错）}不要把 N-S 方程说成"只适用于层流"。准确表述是：
它是\textbf{黏性流体（牛顿流体）的运动微分方程}，\emph{层流}时直接是严格方程，
\emph{湍流}时必须先把速度、压强分解为时均量与脉动量再做时均化，得到雷诺方程。
\end{yicuo}
\end{kaodian}

\begin{kaodian}{2}{N-S 方程的几种解析解（简介）}
只有流动\emph{极简单}的情形才能求出解析解，考法限于"认得名字＋会套式子"。

\begin{gongshi}{两平板间的三种经典解（8.1）}
\[ \text{①纯剪切流（库埃特流）：}\ u=\frac{u_0}{h}y,\qquad\qquad
   \text{②泊肃叶流：}\ u=\frac{1}{2\mu}\frac{\dd p}{\dd x}\left(y^2-hy\right) \]
\[ \text{③库特流（叠加）：}\ u=\frac{u_0}{h}y+\frac{1}{2\mu}\frac{\dd p}{\dd x}\left(y^2-hy\right) \]
\textbf{含义：}①上下板间由板速带动的线性速度分布（无压强梯度）；
②两板均不动、由流向压强梯度 $\dd p/\dd x$ 驱动的抛物线速度分布；
③上板以 $u_0$ 匀速运动\emph{同时}存在压强梯度的库特流，其速度＝①与②的\textbf{叠加}。\\
\textbf{教材例题（例 8.1，书 p156--157）：}这是教材 8.1 节\textbf{唯一}的精确解例题——教材原文：
"设黏性流体在两无限长的水平平板间作恒定层流流动，上板移动速度为 $U_1$，下板移动速度为 $U_2$，
如图 8.3 所示。已知两板间距为 $2h$，质量力可忽略不计，试求两平板间的速度分布。"
教材所用边界条件为 $y=h$ 时 $u=U_1$、$y=-h$ 时 $u=U_2$；所得速度分布同时含\emph{板速项}与\emph{压强梯度项}，
教材原话："如果两平板固定不动，$u=\cdots$"，即退化为②的抛物线（泊肃叶）分布。
圆管泊肃叶流同属 N-S 方程的精确解，教材放在\textbf{第 5 章}（本讲义第 5 章已讲，此处不重复）。\\
\textbf{适用条件：}两平行板、层流、不可压缩牛顿流体、充分发展（$\partial u/\partial x=0$）、定常。
\end{gongshi}

\begin{tishi}
\textbf{订正（原表述 $\to$ 新表述）：}原稿把"\emph{同心圆环间的层流}（轴向压力驱动，速度分布为含 $\ln r$ 的形式）"
写成了教材 8.1 节的精确解之一。核对教材：\textbf{8.1 节并未给出同心圆环间的层流解}，
该节在"下面介绍求解精确解的例题"之后\textbf{只给出例 8.1（两无限长水平平板间恒定层流）}（书 p156--157），
教材原话为"N-S 方程的精确解，虽然为数不多……下面介绍求解精确解的例题"。
新表述：本节的三种平板解对应\textbf{教材例 8.1}；圆管泊肃叶流见\textbf{教材第 5 章}；
同心圆环间的层流解\textbf{教材本章未列}，本讲义仅作扩展提及，\emph{不标教材式号}。
\end{tishi}

\begin{tishi}
\textbf{"为什么只有这几个解"（很可能考的一句话）：}能求出解析解的情形，
都是\emph{流场几何简单到使惯性项 $(\bv\cdot\grad)\bv$ 恒为零或可忽略}的情形
（如平行流动中 $u_y=u_z=0$、$\partial u/\partial x=0$，迁移加速度整项消失）。
一旦惯性项真正起作用，方程就变成非线性的，解析解一般不存在。
\end{tishi}

\textbf{常考题型：}填空/选择（"两平板间由压强梯度驱动的层流速度分布是\underline{\ 抛物线\ }"）；
判断（"库特流的速度是纯剪切流与泊肃叶流的叠加"）。
\textbf{重要度 \xstarfull{2}}\quad\yiju{E2"第六到第八章只有少量填空/判断/选择，掌握基本概念＋公式套用即可"；E6 教材 8.1 节}\quad\nandu{易}

\begin{banfa}
判型三看：\textbf{看板动不动}（动$\to$剪切项$\frac{u_0}{h}y$）、\textbf{看有没有 $\dd p/\dd x$}（有$\to$抛物项）、
\textbf{看两者是否同时有}（同时有$\to$叠加成库特流）。三种情况用\emph{同一个}表达式就能覆盖。
\end{banfa}
\end{kaodian}

\section{边界层的基本概念}

\begin{kaodian}{3}{边界层的定义与三重含义}
\textbf{是什么（教材口径）：}当\textbf{雷诺数很大}时，黏性作用只集中在\textbf{紧贴壁面的极薄一层}内，
这一层称为\textbf{边界层}（附面层）；边界层以外的主流区速度梯度很小，黏性作用可以忽略，
可按\emph{理想流体}处理。

\textbf{教材原文（书 p157，8.2.1 首段，务必按原话背）：}"黏性流体流经固体时，固体边界上的流体质点黏附在固体表面边界上，
与边界没有相对运动，称为\textbf{无滑移条件}。在固体边界的外法线方向上流速从零迅速增大，
在边界附近的流区存在着相当大的速度梯度。在这个流区内黏性作用不能忽略，
边界附近的这个流区就称为\textbf{边界层}（或\textbf{附面层}）。边界层以外的流区，黏性的作用可以略去，可看作是理想流体。"

\textbf{教材的"两个流动区域"模型（教材原话，书 p157）：}"这样，就将大雷诺数流动情况视为由两个性质不同的流动所组成：
一是固体边界附近的边界层流动（如图 8.4 的流动区域 1），黏性作用不能忽略；
另一个是边界层以外的流动（图 8.4 的流动区域 2），按势流理论来求解，
而边界层内的流动，以 N-S 方程为依据，根据问题的物理特点，给予简化处理来求解。"

\textbf{边界层理论的提出者与年份（教材原话，书 p157）：}"1904 年普朗特对此进行研究，结合实验，提出了\textbf{边界层理论}，
对解决大雷诺流动问题提供了求解方法。"
\textbf{教材补一句（书 p157）：}"黏性流体的运动微分方程（N-S 方程），目前只有对最简单边界条件下的少数问题才能求得精确解。
如对于小雷诺数情况，可以略去全部惯性力项，得到简化的线性方程，求得近似解。但是，在实际工程中，大多数是大雷诺数情况，
求解很困难，所以必须寻找新的方法。"

\textbf{边界层的三重含义（概念题就考这三条）：}
\begin{center}
\begin{tabularx}{\textwidth}{@{}lX@{}}
\toprule
含义 & 说明 \\
\midrule
① 一层\textbf{薄}层 & 厚度 $\delta$ 远小于物体特征长度 $l$，即 $\delta\ll l$；且 $\delta/l\sim 1/\sqrt{\Rey_l}$ \\
② \textbf{速度梯度大} & 层内 $\dfrac{\partial u}{\partial y}\sim\dfrac{U}{\delta}$，比流向梯度 $\dfrac{U}{l}$ 大 $\sqrt{\Rey_l}$ 倍，是黏性切应力显著的原因 \\
③ 流动\textbf{有旋} & 层内存在大速度梯度与旋涡，是有旋流动；层外主流可视为无旋的势流 \\
\bottomrule
\end{tabularx}
\end{center}

\textbf{边界层概念的物理图像：}物体前缘处流体与壁面之间的无滑移条件使壁面附近速度陡降；
沿程不断把动量通过黏性向内传递，因此边界层\textbf{越来越厚}。
典型数值（平板层流）：$\delta/x=5.0/\sqrt{\Rey_x}$，例如 $\Rey_x=4\times10^5$ 时 $\delta\approx0.008x$，
即只有板长千分之八——但它是\emph{相对}薄，不是绝对薄。

\textbf{常考题型：}填空（"边界层是指\underline{\ 紧贴壁面、黏性作用显著的薄层\ }"）；
简答（"边界层的含义/为什么引入边界层概念"）；
选择（"边界层外的流动按（理想流体／黏性流体）处理"$\to$理想流体）。
\textbf{重要度 \xstarfull{3}}\quad\yiju{E1 slide16 填空考"名词定义"、简答考"定义/分类/特点"；E1 slide33 第八章＝"选择填空题考点"；E2"掌握基本概念"}\quad\nandu{易}

\begin{banfa}
\textbf{"三字经"背法：}边界层＝\textbf{薄、陡、旋}（薄层、梯度陡、有旋）。
答题时先写"壁面附近 $\delta\ll l$ 的薄层"，再补"层内 $\partial u/\partial y$ 极大、黏性力与惯性力同量级"，最后加"层外按势流处理"——三句话拿满分。
\end{banfa}

\begin{yicuo}
\textbf{易错 1：}"薄"是\textbf{相对}的。$\delta/l\sim1/\sqrt{\Rey_l}$ 随 $\Rey$ \emph{增大而减小}，
但 $\delta$ 的\emph{绝对值}沿程仍然增大。飞机机翼边界层可以有几厘米厚，不能想当然认为它"很薄所以可以忽略不计"。
\end{yicuo}
\end{kaodian}

\begin{kaodian}{3}{边界层厚度 $\delta$ 的定义与沿程增长}
\begin{gongshi}{边界层厚度的工程定义（8.2）}
\[ \delta(x):\quad u(x,y)\big|_{y=\delta}=0.99\,U_\infty \]
\[ \text{量级关系：}\quad \frac{\delta}{l}\sim\frac{1}{\sqrt{\Rey_l}}\quad\text{或}\quad \delta\propto\sqrt{\frac{\nu x}{U}} \]
\textbf{含义：}把\emph{速度达到来流速度（外部势流速度）$99\%$} 的位置规定为边界层外缘，
该处到壁面的法向距离即边界层厚度 $\delta$。用 $99\%$ 而非 $100\%$ 是因为从层内到层外的过渡是\emph{渐近}的，
不存在一个速度恰好等于 $U_\infty$ 的明确界面，所以 $\delta$ 是一个\textbf{人为约定的工程量}。\\
\textbf{教材原话（书 p158，8.2.2）：}"一般规定 $u=0.99U_0$ 的地方可看作是边界层外边缘，
可以认为边界层厚度（\emph{几何厚度}）是沿固体表面外法线方向从 $u=0$ 到 $u=0.99U_0$ 的一段距离。"
教材并强调\textbf{边界层厚度是 $x$ 的函数}。\\
\textbf{教材式号：}量级关系 $\delta\propto\sqrt{\nu x/U_0}$ 即（教材式 8.16）（书 p161）；
平板层流的 $\delta=5.0x/\sqrt{\Rey_x}$ 即（教材式 8.35）（书 p165）。\\
\textbf{适用条件：}适用于大雷诺数下壁面附近的流动；$\delta$ 依赖约定（若改用 $0.995U_\infty$ 则 $\delta$ 略大），
故涉及 $\delta$ 的一切公式都必须"同一约定配套使用"。
\end{gongshi}

\textbf{沿程增长的物理原因：}流体从前缘开始，壁面处速度为零，靠黏性把外部动量逐层向内"扩散"；
扩散所需时间随流向距离增加，故 $\delta$ 随 $x$ 增大，且 $\delta\propto\sqrt{x}$（层流）。
\textbf{平板}上 $U_\infty$ 不变时 $\delta=5.0x/\sqrt{\Rey_x}$；层流时 $\delta$ 是抛物线型的缓增，
转捩成湍流后 $\delta$ 增长明显加快（$\delta\propto x^{4/5}$）。

\textbf{常考题型：}填空（"边界层厚度定义为 $u=$ \underline{\ } $U_\infty$ 处的 $y$ 值"$\to$0.99）；
选择（"边界层厚度沿程如何变化"$\to$增大）；
小计算（给出 $\Rey_x$ 求 $\delta/x$）。
\textbf{重要度 \xstarfull{3}}\quad\yiju{E1 slide16 填空考"名词定义与小计算"；E2"掌握基本概念＋公式套用"；E6 教材 8.2 节}\quad\nandu{易}

\begin{banfa}
\textbf{约定的数字只有两个要记：}厚度用 $\mathbf{0.99}$，转捩用 $\mathbf{5\times10^5}$。
\textbf{顺手自检：}$\delta/x=5.0/\sqrt{\Rey_x}$ 中，$\Rey_x=10^6$ 时 $\delta/x=0.5\%$；
考试里若算出 $\delta/x$ 达到百分之几十，一定是把 $\Rey_x$ 算错了。
\end{banfa}
\end{kaodian}

\begin{kaodian}{3}{边界层内的两种流态与转捩点}
\textbf{是什么：}边界层内部同样存在\textbf{层流边界层}与\textbf{湍流边界层}两种流态；
由层流转为湍流的现象称为\textbf{转捩}，转捩发生的位置称为\textbf{转捩点}。
平板绕流中，前缘开始的层流边界层经过一段距离后转捩为湍流。

\begin{gongshi}{边界层流态的判据（8.2）}
\[ \Rey_x=\frac{U_\infty x}{\nu},\qquad\qquad \Rey_{x,c}\approx5\times10^5 \]
\textbf{含义：}以\emph{距前缘的距离 $x$ 为特征长度}构造的局部雷诺数。
$\Rey_x<5\times10^5$ 时为层流边界层，$\Rey_x>5\times10^5$ 后转捩为湍流边界层。\\
\textbf{教材式号：}$\Rey_x=U_0x/\nu$ 即（教材式 8.11）（书 p159）。\\
\textbf{教材原话（书 p159）：}"沿流动方向，随着 $x$ 增加，雷诺数增大，当其达到一定数值后，
边界层内流动经过一过渡段后转变为湍流，成为湍流边界层。由层流边界层转变为湍流边界层的点 $x_c$ 设为\textbf{转捩点}，
对应的雷诺数称为\textbf{临界雷诺数} $\Rey_{x,cr}$。对于光滑平板而言，$\Rey_{x,cr}$ 的范围为 $3\times10^5\sim3\times10^6$，
一般取 $5\times10^5$。在湍流边界层中，紧贴平板表面亦有一层极薄的\textbf{黏性底层}。"
\textbf{教材补充（书 p159，边界层概念也适用于管流与明渠）：}"由于受壁面阻滞的影响，靠近管壁或渠壁的流体在进口附近形成边界层，
其厚度 $\delta$ 随离进口的距离的增加而加大。当边界层发展到管轴或渠道自由表面后，流体的运动都处于边界层内，
此后流速分布不再变化，形成均匀流动。从进口发展到均匀流的长度，称为\textbf{进口段长度}，用 $L'$ 表示。
对于圆管层流，$L'=0.065\Rey d$；对于圆管湍流，$L'=(50\sim100)d$。"\\
\textbf{适用条件：}适用于光滑平板、来流湍流度较小的情形；
转捩点位置受来流湍流度、壁面粗糙度、压力梯度与壁面温度影响很大，
实际可在 $3\times10^5\sim3\times10^6$ 之间变动，故 $5\times10^5$ 只是\textbf{工程参考值}（教材口径与本讲义完全一致）。
\end{gongshi}

\textbf{湍流边界层的分层（必背）：}紧贴壁面是极薄的\textbf{黏性底层}（层流），
再往外是\textbf{过渡区}，最外层是\textbf{湍流核心区}。
湍流边界层内部由于脉动掺混，速度剖面比层流\emph{饱满}（同一 $\delta$ 下，近壁速度更大、
壁面切应力也更大），故湍流边界层的\emph{摩擦阻力更大}，
但\emph{抵抗逆压梯度的能力更强}、\emph{不易分离}。

\textbf{常考题型：}选择（"平板边界层由层流转为湍流的临界雷诺数约为"$\to5\times10^5$）；
判断（"边界层一定是层流"$\to$错）；填空（"湍流边界层由黏性底层、过渡区和\underline{\ 湍流核心区\ }组成"）。
\textbf{重要度 \xstarfull{3}}\quad\yiju{E1 slide16 选择/判断考"概念辨析"、填空考"名词定义"；E2"掌握基本概念"；E6 教材 8.2 节}\quad\nandu{易}

\begin{banfa}
\textbf{"两个 $\Rey$ 分清楚"对照表（考点即在此）：}
\begin{center}
\begin{tabularx}{\textwidth}{@{}lXX@{}}
\toprule
 & 管内流动（第 5 章） & 平板边界层（本章） \\
\midrule
特征长度 & 管径 $d$ & 距前缘距离 $x$ \\
判据 & $\Rey_d=\dfrac{vd}{\nu}$ & $\Rey_x=\dfrac{U_\infty x}{\nu}$ \\
临界值 & $2000$（工程取值） & $5\times10^5$ \\
\bottomrule
\end{tabularx}
\end{center}
题目里出现"$d$"就用第一个，出现"$x$、距前缘"就用第二个，\textbf{绝不能互换}。
\end{banfa}

\begin{yicuo}
\textbf{易错 3：}边界层的层流/湍流与\emph{管内的}层流/湍流是\textbf{两套判据}：
管内按 $\Rey_d$ 且临界值是 $2000$；平板边界层按 $\Rey_x$ 且临界值是 $5\times10^5$。
把 $2000$ 套到平板上（或反之）是最典型的失分点。
\end{yicuo}
\end{kaodian}

\begin{kaodian}{3}{边界层内的压强分布}
\begin{gongshi}{边界层内沿法向压强不变（8.2、8.3，核心结论）}
\[ \frac{\partial p}{\partial y}\approx0\quad\Longrightarrow\quad p=p(x) \]
\textbf{含义：}边界层很薄，法向坐标 $y$ 方向上的惯性效应（离心效应）可忽略，
故\textbf{压强沿壁面法向不变}，边界层内任意一点的压强\emph{都等于}该处边界层外缘的压强；
$(x,y)$ 处的二维压强场退化为一维的 $p=p(x)$。这就是"边界层内压强由外部势流决定"的准确表述。\\
\textbf{教材原话（书 p160--161，数量级分析的结论）：}"压强 $p$ 在 $y$ 方向的变化非常小，认为 $p$ 仅随 $x$ 改变，
即 $p=p(x)$。说明整个边界层厚度方向压强不变，相同 $x$ 坐标边界层内、外压强相等。"
\textbf{教材式号：}该结论来自（教材式 8.12）平面恒定边界层的连续性方程与 N-S 方程、及（教材式 8.13）
$\partial u_x/\partial x=-\partial u_y/\partial y$（数量级分析用）（书 p159--160）。\\
\textbf{适用条件：}壁面曲率半径远大于边界层厚度（薄层近似）且流动未分离。
圆柱/机翼等曲面壁面只要 $\delta\ll$ 曲率半径仍可用；
但\textbf{分离区、尾流区、激波区}不成立。
\end{gongshi}

\textbf{由外部势流确定压强梯度（解题必经的一步）：}
边界层外缘是理想流体，沿流线满足伯努利方程 $\dfrac{p}{\rho}+\dfrac{U^2}{2}=$常数，对 $x$ 求导得
\[ \frac{\dd p}{\dd x}=-\rho U\frac{\dd U}{\dd x} \]
\textbf{讨论三种情形：}
\begin{center}
\begin{tabularx}{\textwidth}{@{}lXX@{}}
\toprule
情形 & 判据 & 边界层表现 \\
\midrule
顺压梯度 & $\dfrac{\dd p}{\dd x}<0$（$\dd U/\dd x>0$） & 压强沿程下降，流体被加速，剖面更饱满，\textbf{不易分离} \\
零压梯度 & $\dfrac{\dd p}{\dd x}=0$（平板 $U=$常数） & 平板绕流的理想情形 \\
逆压梯度 & $\dfrac{\dd p}{\dd x}>0$（$\dd U/\dd x<0$） & 压强沿程升高，流体减速，\textbf{可能分离}（见 8.7） \\
\bottomrule
\end{tabularx}
\end{center}

\textbf{常考题型：}填空（"边界层内压强沿法向\underline{\ 不变\ }，即 $p=p(x)$"）；
选择（"边界层内压强梯度 $\dd p/\dd x$ 由（外部势流／边界层内的黏性／壁面切应力）决定"$\to$外部势流）；
简答（"为什么边界层内压强只沿流向变化"）。
\textbf{重要度 \xstarfull{3}}\quad\yiju{E1 slide16 选择/判断考"概念辨析"、填空考"名词定义"；E2"掌握基本概念"；E6 教材 8.2、8.3 节\quad}\nandu{中}

\begin{banfa}
\textbf{这是"白送分"的突破口：}题目只要给了物体表面的\emph{压强分布}或\emph{外部势流速度分布}，
就能直接用 $\dfrac{\dd p}{\dd x}=-\rho U\dfrac{\dd U}{\dd x}$ 换算，把边界层方程右边的压强项写成 $U$ 的函数。
\textbf{一句话答题模板：}"因边界层极薄，$\partial p/\partial y\approx0$，层内压强沿法向不变、
等于层外势流压强，故由伯努利方程得 $\dfrac{\dd p}{\dd x}=-\rho U\dfrac{\dd U}{\dd x}$。"
\end{banfa}

\begin{yicuo}
\textbf{易错 2：}不要以为"边界层内压强也跟着黏性耗散掉很多"。恰恰相反——
边界层内的\textbf{黏性只在动量方程里起作用，在压强上不产生法向落差}。
压强是\emph{外部势流说了算}的已知量，这是边界层方程能被简化为抛物型方程的关键前提。
\end{yicuo}
\end{kaodian}

\begin{kaodian}{2}{边界层概念的适用条件}
\textbf{成立条件（三条）：}
\begin{enumerate}
  \item \textbf{大雷诺数}：$\Rey$ 足够大，$\delta/l\sim1/\sqrt{\Rey}$ 才足够小，"薄层＋层外无黏"才成立；
  低 $\Rey$（如 $\Rey<10$ 的绕流球）时黏性遍布全流场，边界层概念\textbf{不适用}；
  \item \textbf{壁面附近、未分离}：只描述贴壁薄层与主流的衔接；\emph{分离点以后}边界层方程失效；
  \item \textbf{流体连续}：连续介质假设成立（稀薄气体、激波后强间断不适用）。
\end{enumerate}

\textbf{不适用的区域（选择/判断题的陷阱材料）：}
\begin{center}
\begin{tabularx}{\textwidth}{@{}lX@{}}
\toprule
区域 & 为什么不适用 \\
\midrule
尾流区 & 物体后面黏性影响扩散到整个尾迹，不再"薄"，必须用完整 N-S 方程或湍流模型 \\
分离区 & $\partial u/\partial y|_{y=0}=0$ 之后出现回流，方程类型与边界条件全部改变 \\
激波区 & 存在强间断、压缩性显著，连续介质与不可压缩前提被破坏 \\
低 $\Rey$ 绕流 & 黏性作用遍布全流场，没有"层外势流"可言 \\
\bottomrule
\end{tabularx}
\end{center}

\textbf{常考题型：}判断（"边界层理论适用于任何雷诺数的绕流"$\to$错）；
选择（"下列哪个区域不能用边界层理论"$\to$尾流/分离区）。
\textbf{重要度 \xstarfull{2}}\quad\yiju{E2"第六到第八章只有少量填空/判断/选择"；E1 slide16 判断考"概念辨析"；E6 教材 8.2 节}\quad\nandu{易}

\begin{banfa}
\textbf{"四不适用"口诀：}尾流、分离、激波、低 $\Rey$。
凡是题目描述里出现"物体后面很远的地方""回流""激波""$\Rey$ 很小"，就要立刻想到边界层理论已失效。
\end{banfa}
\end{kaodian}

\section{边界层方程组及边界条件}

\begin{kaodian}{3}{量级比较与普朗特边界层方程}
\textbf{思路（简答常考）：}把完整的不可压缩 N-S 方程在\emph{边界层薄层}内做\textbf{量级比较}，
把量级远小于惯性项的黏性项与法向压强项丢掉，得到\textbf{普朗特边界层方程}。

\textbf{量级比较的依据与结果（必背表格）：}设流向特征长度 $l$、法向厚度 $\delta$、外部速度 $U$。
\begin{center}
\begin{tabularx}{\textwidth}{@{}lXc@{}}
\toprule
量 & 依据 & 量级 \\
\midrule
$x$ & 顺流方向变化缓慢 & $l$ \\
$y$ & 只跨过薄薄一层 & $\delta$ \\
$u$ & 与外部势流同级 & $U$ \\
$v$ & 由连续性方程 $\dfrac{\partial u}{\partial x}+\dfrac{\partial v}{\partial y}=0$ 反推 &
$\dfrac{U\delta}{l}$ \\
$\dfrac{\partial^2 u}{\partial y^2}$ & 黏性与惯性同量级 & $\dfrac{U}{\delta^2}$ \\
\bottomrule
\end{tabularx}
\end{center}

由"惯性项 $\sim$ 黏性项"得 $U^2/l\sim\nu U/\delta^2$，即
\[ \delta\sim\sqrt{\frac{\nu l}{U}}\ ,\qquad \frac{\delta}{l}\sim\frac{1}{\sqrt{\Rey_l}} \]
这同时\emph{证明了}边界层厚度与 $\sqrt{x}$ 成正比，也说明 $v\ll u$（$v/u\sim\delta/l$）。

\textbf{量级比较的结论（丢掉哪些项）：}
\begin{center}
\begin{tabularx}{\textwidth}{@{}lX@{}}
\toprule
方程 & 简化结果 \\
\midrule
连续方程 & 形式\textbf{不变}：$\dfrac{\partial u}{\partial x}+\dfrac{\partial v}{\partial y}=0$ \\
$x$ 向动量方程 & 黏性项只保留 $\nu\dfrac{\partial^2 u}{\partial y^2}$（$\dfrac{\partial^2 u}{\partial x^2}$ 小 $\Rey_l$ 倍，丢弃） \\
$y$ 向动量方程 & 各黏性项与惯性项均比 $x$ 向小，只剩 $\dfrac{\partial p}{\partial y}=0$，即 $p=p(x)$ \\
\bottomrule
\end{tabularx}
\end{center}

\begin{gongshi}{普朗特边界层方程（8.3，必背）}
\[ \text{连续：}\ \frac{\partial u}{\partial x}+\frac{\partial v}{\partial y}=0 \]
\[ \text{动量：}\ u\frac{\partial u}{\partial x}+v\frac{\partial u}{\partial y}
   =-\frac{1}{\rho}\frac{\dd p}{\dd x}+\nu\frac{\partial^2 u}{\partial y^2} \]
\[ \text{压强项：}\ \frac{\dd p}{\dd x}=-\rho U\frac{\dd U}{\dd x}\quad(\text{由外部势流伯努利方程}) \]
\textbf{含义：}在薄层近似下，二维不可压缩定常层流边界层的基本方程组。
它与 N-S 方程的本质区别是：\textbf{法向动量方程被压强方程取代}，$p$ 由外部势流给出；
方程对 $x$ 是\emph{抛物型}的，可以从前缘开始沿 $x$ \emph{逐步向前推进}求解\\
（而 N-S 方程是椭圆型的，必须全场耦合求解）——这正是边界层理论把问题"化难为易"的关键。\\
\textbf{教材式号：}\textbf{（教材式 8.14）普朗特边界层方程}；另有（教材式 8.15）边界层方程组的另一种形式
（压强项换成 $U_0$ 形式）（书 p161）。
\textbf{教材原话（书 p161）：}"普朗特边界层方程较 N-S 方程是大大简化了……方程由\textbf{椭圆方程}变成了\textbf{抛物线方程}。
问题的求解域由一个二维的无穷域变成了一个半无限的长条域。"（与本讲义下文"抛物型可逐步推进求解"完全对应）\\
\textbf{适用条件：}大 $\Rey$、薄边界层、\emph{未分离}、二维（或轴对称）层流、
不可压缩、物性为常数、壁面曲率半径远大于 $\delta$。
\end{gongshi}

\textbf{常考题型：}填空（"边界层方程中 $y$ 向动量方程简化为\underline{\ $\partial p/\partial y=0$\ }"）；
选择（"边界层方程中保留的黏性项是"$\to$ $\nu\partial^2u/\partial y^2$）；
简答（"写出边界层方程并说明各项含义"）。
\textbf{重要度 \xstarfull{3}}\quad\yiju{E1 slide33 第八章＝"选择填空题考点"；E1 slide16 选择考"概念辨析"；E2"掌握基本概念与公式套用"；E6 教材 8.3 节}\quad\nandu{中}

\begin{banfa}
\textbf{"只留一个二阶项"记忆法：}N-S 的黏性项有三个 $\partial^2/\partial x^2$、$\partial^2/\partial y^2$、$\partial^2/\partial z^2$；
边界层里\textbf{只有跨过薄层的 $\partial^2/\partial y^2$ 幸存}，其余丢进"小一个 $\Rey_l$ 量级"里。
考试写方程时若多写了 $\nu\partial^2u/\partial x^2$，就会被扣"未做量级比较"的分。
\end{banfa}
\end{kaodian}

\begin{kaodian}{3}{边界条件}
\begin{gongshi}{边界层方程的边界条件（8.3）}
\[ y=0:\quad u=0,\quad v=0\qquad(\text{无滑移、无穿透}) \]
\[ y=\delta\ (\text{或 }y\to\infty):\quad u\to U(x)\qquad(\text{与外缘势流衔接}) \]
\[ \text{压强：}\quad \frac{\dd p}{\dd x}=-\rho U\frac{\dd U}{\dd x} \]
\[ \text{前缘条件：}\quad x=0:\ u=U_\infty\ (\text{均匀来流}) \]
\textbf{含义：}① \textbf{壁面无滑移}：黏性流体紧贴壁面的流体质点与壁面\emph{没有相对速度}，
故 $u=v=0$；这是边界层（乃至一切黏性流动）产生的根源。
② \textbf{外缘衔接}：边界层外缘速度等于外部势流速度 $U(x)$，
由外部势流解（势流理论，见第 7 章）提供；
③ 压强梯度不再是独立未知量，由外部势流经伯努利方程换算得到。\\
\textbf{教材式号：}（教材式 8.17）恒定不可压二维边界层的控制方程及边界条件（书 p162）。
\textbf{教材原话（书 p161）：}"当 $y=0$ 时（壁面处），$u_x(x,0)=u_y(x,0)=0$；当 $y=\delta$ 或 $y\to\infty$ 时，$u_x(x,y)=U_0(x)$。"
\textbf{教材另一句（书 p162，说明流函数的作用）：}"引入流函数后，可以用一个流函数 $\psi(x,y)$ 代替两个速度分量
$u_x(x,y)$、$u_y(x,y)$，使得两个偏微分方程合并为一个。"\\
\textbf{适用条件：}固壁、连续介质、未分离；壁面运动时无滑移条件改为 $u=U_{\text{wall}}$。
\end{gongshi}

\textbf{为什么必须有无滑移条件：}若允许壁面处流体滑移，则壁面处 $\partial u/\partial y$ 不再是 $\sim U/\delta$，
边界层的整套量级比较和黏性效应都失去来源。无滑移是"边界层存在"的逻辑前提，
也是\emph{理想流体理论无法给出阻力}的根源——理想流体允许在壁面滑移，算不出摩擦阻力。

\textbf{常考题型：}填空（"边界层方程的壁面边界条件是\underline{\ $u=v=0$\ }"）；
简答（"写出边界层方程及其边界条件"）；判断（"黏性流体在壁面处可以滑移"$\to$错）。
\textbf{重要度 \xstarfull{3}}\quad\yiju{E1 slide16 填空考"名词定义"、简答考"定义/分类/特点"；E2"掌握基本概念"；E6 教材 8.3 节}\quad\nandu{易}

\begin{banfa}
\textbf{边界条件"三层检验"：}拿到边界层题目，先写\emph{壁面两条}（$u=0$、$v=0$），
再写\emph{外缘一条}（$u=U(x)$），最后写\emph{压强换算一条}（$\dd p/\dd x=-\rho U\,\dd U/\dd x$）。
四行写完，即使没往下算，简答题也已经拿到了大部分分数。
\end{banfa}
\end{kaodian}

\section{边界层方程的相似性解}

\begin{kaodian}{2}{布拉修斯解（平板层流边界层）}
\textbf{思路（相似性解的要义）：}平板边界层外部 $U_\infty=$常数，是\emph{零压梯度}情形。
布拉修斯提出：不同 $x$ 断面上的速度剖面\emph{彼此相似}，可用一个组合变量把它们"归一化"。
引入无量纲法向坐标与流函数
\[ \eta=y\sqrt{\frac{U_\infty}{\nu x}}\ ,\qquad \psi=\sqrt{\nu xU_\infty}\,f(\eta)
   \ \Longrightarrow\ u=U_\infty f'(\eta) \]
代入普朗特方程后，偏微分方程化为\textbf{常微分方程}
\[ 2f'''+ff''=0 \]
边界条件：$\eta=0$ 处 $f=f'=0$（无滑移）；$\eta\to\infty$ 处 $f'\to1$（外缘衔接）。
数值求得 $f''(0)=0.332$。\textbf{这就是"相似性解"：把 $x$、$y$ 两个自变量合成一个 $\eta$，
使不同断面的速度剖面成为同一条曲线。}

\textbf{教材式号（8.4 节全套，书 p162--164）：}（教材式 8.18）$u(x,y)/U_0(x)=f(\eta)$；
（教材式 8.19）引入流函数后的方程与边界条件；（教材式 8.20）$\psi(x,y)=U_0(x)g(x)f(\eta)$；
（教材式 8.21）$f'''+\alpha ff''+\beta(1-f'^2)=0$；（教材式 8.22）$\alpha$、$\beta$ 的定义；
（教材式 8.23）$(2\alpha-\beta)\nu x=U_0g^2$；（教材式 8.24）$U_0(x)=Cx^m$；（教材式 8.25）$g(x)$ 的形式；
（教材式 8.26）$2f'''+ff''=0$；（教材式 8.27）$g(x)=\sqrt{\nu x/U_0}$；（教材式 8.28）$\eta=y\sqrt{U_0/(\nu x)}$；
（教材式 8.29）$\psi=\sqrt{\nu xU_0}f(\eta)$；（教材式 8.30）$u=U_0f'(\eta)$；
（教材式 8.31）$u_y=\dfrac12\sqrt{\nu U_0/x}\,(\eta f'-f)$。
\textbf{教材表：}\textbf{表 8.1}"平板边界层的解的结果"（书 p164--165）给出 $f(\eta)$、$f'(\eta)$、$f''(\eta)$ 的数值表，
查表得 $\eta=5.0$ 时 $f'\approx0.99$（教材据此认定此处即边界层外缘）、$\eta=0$ 时 $f''(0)=0.332$。

\begin{tishi}
相似性解的意义有三条：① 把偏微分方程降为常微分方程，可数值求解；
② 说明层流边界层速度剖面\textbf{自相似}（去量纲后与 $x$ 无关）；
③ 由此得到下面一组\textbf{可直接套用的公式}。\emph{推导过程不需背，结果必须背。}
\end{tishi}

\begin{tishi}
\textbf{教材补全（相似性解的存在条件，书 p163，教材原话）：}"只有当边界层外部势流速度为\textbf{幂函数形式}时，
边界层方程才有相似性解，边界层偏微分方程才能转化为常微分方程。"即（教材式 8.24）$U_0(x)=Cx^m$；
上式中指数 $m=\beta/(2\alpha-\beta)$，$C$ 为常数；"式（8.21）中的系数 $\beta=0$ 表示流体顺流绕过平板形成平板边界层流动"。
\textbf{订正/补全：}本讲义原只说"布拉修斯提出不同 $x$ 断面剖面相似"，未说明相似性解\emph{存在条件}——
按教材补出：\textbf{只有当外部势流速度是幂函数 $U_0=Cx^m$ 时才有相似性解}，平板（$m=0$、$\beta=0$）只是其中最简单的一例。
\textbf{教材自身印误（两说并列，不删原说法）：}教材原文写"边界层方程的相似性解最初由德国科学家布拉修斯进行了研究，
\textbf{1980 年}他在其博士论文中详细讨论了这个问题，这是第一个应用普朗特边界层理论的具体例子。"（书 p163）
\emph{另一说}：流体力学史上布拉修斯（H.~Blasius）的博士论文为 \textbf{1908 年}（普朗特提出边界层理论为 1904 年），
"1980 年"应为教材误印或 OCR 误读。\textbf{本讲义两说并列}：正文只写"布拉修斯在其博士论文中首次给出该解"，不写年份；
考场若照教材原文答"1980 年"亦可，若答"1908 年"亦有据——切忌自造其它年份。
另：教材表 8.1 表头所载解的来源人名 OCR 作"\textbf{豪华斯}"（通行译名为"\emph{豪沃斯}"，即 Howarth），
此处存疑待核，本讲义不据此立论。
\end{tishi}

\textbf{常考题型：}选择/填空（"布拉修斯解适用于（平板／零压梯度）层流边界层"）；
简答（"什么是边界层方程的相似性解"）。
\textbf{重要度 \xstarfull{2}}\quad\yiju{E2"第六到第八章只有少量填空/判断/选择，掌握基本概念＋公式套用即可"；E6 教材 8.4 节}\quad\nandu{难}

\begin{banfa}
\textbf{认标志就能选对：}出现"$\eta=y\sqrt{U/(\nu x)}$""$f''(0)=0.332$""平板层流"三者之一，
就是布拉修斯解。它\textbf{只能用于平板}（$U=$常数、$\dd p/\dd x=0$），
一旦给出曲面或非零压强梯度，就不能套这组数。
\end{banfa}
\end{kaodian}

\begin{kaodian}{2}{平板层流边界层的厚度、切应力与阻力公式}
\begin{gongshi}{布拉修斯解的四个结果（8.4，必背）}
\[ \delta=\frac{5.0x}{\sqrt{\Rey_x}}\ ,\qquad\qquad
   \tau_w=0.332\frac{\rho U_\infty^2}{\sqrt{\Rey_x}} \]
\[ C_f=\frac{\tau_w}{\frac{1}{2}\rho U_\infty^2}=\frac{0.664}{\sqrt{\Rey_x}}\ ,\qquad\qquad
   C_D=\frac{1.328}{\sqrt{\Rey_l}} \]
\[ F_d=\frac{1}{2}\rho U_\infty^2\,bl\,C_D=0.664\,b\sqrt{\mu\rho l}\,U_\infty^{1.5} \]
\textbf{含义：}$\delta$＝边界层厚度（$u=0.99U_\infty$ 处）；
$\tau_w$＝壁面局部切应力；$C_f$＝\textbf{局部摩擦阻力系数}（该 $x$ 处单位面积摩擦力与 $\frac12\rho U^2$ 之比）；
$C_D$＝\textbf{整板平均摩擦阻力系数}；$F_d$＝宽 $b$、长 $l$ 平板单面的总摩擦阻力。\\
\textbf{教材式号：}$C_f=0.664/\sqrt{\Rey_x}$ 即（教材式 8.32）；一侧摩擦阻力 $D=\int_0^L\tau_0b\,\dd x$ 即（教材式 8.33）；
摩擦阻力系数 $C=1.32/\sqrt{\Rey_L}$ 即（教材式 8.34）；$\delta=5.0x/\sqrt{\Rey_x}$ 即（教材式 8.35）（书 p165）。
排挤厚度 $\delta_1=1.72x/\sqrt{\Rey_x}$、动量损失厚度 $\theta=0.664x/\sqrt{\Rey_x}$ 在教材中\textbf{无式号}（书 p158、p165）。
\textbf{教材例题：}（例 8.2，书 p171）长 5 m、宽 2 m 平板，$U_0=0.1$ m/s、$\nu=1.139\times10^{-6}$ m$^2$/s、$\rho=999.1$ kg/m$^3$，
求距首端 1 m 与 4 m 处边界层厚度及平板一面摩擦阻力（教材用（教材式 8.50）$C_f=1.46/\sqrt{\Rey_L}$，得 $\delta(1\ \text{m})=1.85$ cm、$D=0.11$ N）；
（例 8.3，书 p171--172）2 m$\times$8 m 平板两种放法的摩擦阻力比 $F_1/F_2=0.908$——即 E4 2015 年"平板摩擦阻力比"的\textbf{教材同型题}。\\
\textbf{适用条件：}\emph{平板}、\emph{层流边界层}（$\Rey_l<5\times10^5$）、不可压缩、定常、零压梯度、
来流湍流度很低；$F_d$ 为\emph{单面}阻力（双面板要乘 2）；$C_D=1.328/\sqrt{\Rey_l}$ 中的
$\Rey_l=U_\infty l/\nu$ 用\textbf{板长}作特征长度。
\end{gongshi}

\begin{tishi}
\textbf{教材自身印误（两说并列，不删原说法）：}"整板平均摩擦阻力系数"一项，教材\textbf{（教材式 8.34）印作 $C=1.32/\sqrt{\Rey_L}$}（书 p165），
而按紧邻的（教材式 8.33）$D=\int_0^L\tau_0b\,\dd x=0.332\rho U_0^2b\cdot2\sqrt{\nu L/U_0}$ 与 $C_D=2D/(\rho U_0^2 bL)$ 应为
$1.328/\sqrt{\Rey_L}$（恰为 $2\times0.664$）。\emph{另一说}：1.32 系教材按有效数字取两位所致（教材同页（教材式 8.32）写作 0.664 三位有效数字），
故更可能是\textbf{印刷/取位误差}。\textbf{本讲义取 $C_D=1.328/\sqrt{\Rey_l}$}（与 0.664 自洽、且是历年真题通行取值），
同时保留教材原值 1.32 备查。同理，排挤厚度本讲义写 $1.721$，教材（书 p165）写作 $1.72$——同一量的两种取位。
另外注意：\textbf{例 8.2 与例 8.3 用的 $1.46/\sqrt{\Rey_L}$（教材式 8.50）是"动量积分近似解"}，
比布拉修斯精确解的 $1.328/\sqrt{\Rey_L}$ 大，二者\textbf{不可混用}——判卷时按题目所给教材口径取。
\end{tishi}

\textbf{数量级自检（很好用）：}$C_f$ 随 $x$ 增大而减小（$\propto x^{-1/2}$），
说明\emph{前缘附近单位面积摩擦力最大}、越往后越小；而 $\delta$ 越往后越大。
这两件事方向相反，切忌混答。

\textbf{常考题型：}填空（"平板层流边界层 $\delta/x=$ \underline{\ }"）；
计算（给定 $U$、$l$、$b$、$\nu$ 求 $F_d$ 或求两块板的阻力比）；
选择（"$C_f$ 随 $\Rey_x$ 如何变化"$\to$减小）。
\textbf{重要度 \xstarfull{2}}\quad\yiju{E4 2015 年试题计算题含"平板摩擦阻力比"；E2"掌握基本概念＋公式套用即可"；E6 教材 8.4、8.6 节}\quad\nandu{中}

\begin{banfa}
\textbf{平板阻力计算"五步走"（专治"平板摩擦阻力比"这类题）：}
① 由 $U_\infty$、$l$、$\nu$ 算 $\Rey_l=\dfrac{U_\infty l}{\nu}$，看是否小于 $5\times10^5$ 以判定层流；
② 层流取 $C_D=\dfrac{1.328}{\sqrt{\Rey_l}}$；
③ 单面阻力 $F_d=\dfrac12\rho U_\infty^2\,bl\,C_D$，双面乘 2；
④ 若题目问"两块板的阻力比"，先分别算出两个 $\Rey_l$，再按 $F\propto \sqrt{U^3 l}$（同流体温、同宽）或
直接用 $F=\dfrac12\rho U^2b\sqrt{\nu l U}\times1.328$ 求比，\textbf{不要漏掉任何一个变量}；
⑤ 最后做量纲检查：阻力必须是 N，且与 $b$、$l^{1/2}$、$U^{3/2}$ 成正比。

\textbf{秒答技巧：}层流平板的 $F_d=0.664b\sqrt{\mu\rho l}\,U^{1.5}$，
说明\textbf{层流摩擦阻力与来流速度的 1.5 次方成正比}，而与板长的平方根成正比。
题目若给"速度加倍"，阻力变成 $2^{1.5}=2.83$ 倍——这个结论比一步步算快得多。
\end{banfa}
\end{kaodian}

\section{边界层动量积分方程}

\begin{kaodian}{2}{冯·卡门动量积分方程}
\textbf{思路（推导要点）：}在边界层内取一个\textbf{控制体}（长 $\dd x$、高跨过边界层），
对其应用\textbf{动量定理}：\emph{流出与流入控制体的动量通量之差＝作用在控制体上的合力}。
作用力只有三个来源：壁面切应力 $\tau_w\dd x$（阻力方向）、
两侧压强差产生的压强力、以及外缘带入的动量。
把边界层方程对 $y$ 从 $0$ 积分到 $\delta$ 也能得到同一结果。

\begin{gongshi}{卡门动量积分方程（8.5，必背形式）}
\[ \frac{\tau_w}{\rho U^2}=\frac{\dd \delta_2}{\dd x}
   +2\delta_2\frac{\dd U}{\dd x}+\frac{\delta_2}{U}\frac{\dd U}{\dd x}
   =\frac{\dd \delta_2}{\dd x}+\frac{\delta_2}{U}\left(2+H\right)\frac{\dd U}{\dd x} \]
\textbf{平板（$\dd U/\dd x=0$）简化为最常用的形式：}
\[ \frac{\tau_w}{\rho U_\infty^2}=\frac{\dd \delta_2}{\dd x}\ ,\qquad\text{即}\qquad
   \tau_w=\rho U_\infty^2\frac{\dd \delta_2}{\dd x} \]
\textbf{含义：}把动量守恒沿边界层厚度方向积分，把\emph{二维偏微分方程}降为关于 $x$ 的\emph{常微分方程}；
方程中不再出现 $v$ 与法向坐标，只涉及 $\tau_w$、$U$ 与三个厚度参数。
它是"近似解法"的理论基础：\emph{只要假定一个速度剖面，就能解出 $\delta(x)$ 与 $\tau_w(x)$}。\\
\textbf{适用条件：}二维、不可压缩、定常、\emph{不分离}（要求 $\tau_w>0$）；
既适用于层流，也适用于湍流（湍流时 $\tau_w$ 取时均湍流切应力）。
\end{gongshi}

\textbf{常考题型：}填空（"平板零压梯度时卡门方程为\underline{\ $\tau_w=\rho U^2\,\dd\delta_2/\dd x$\ }"）；
简答（"简述动量积分方程的推导思路与意义"）。
\textbf{重要度 \xstarfull{2}}\quad\yiju{E2"第六到第八章只有少量填空/判断/选择"；E1 slide33 第八章＝"选择填空题考点"；E6 教材 8.5 节}\quad\nandu{难}

\begin{banfa}
\textbf{"积分方程解决什么"的标准答法：}
"它把边界层的\textbf{偏微分方程沿法向积分一次}，化为常微分方程，
从而\textbf{不需要知道精确的速度剖面}，只要假定的剖面满足边界条件，
就能求出工程上够用的 $\delta$、$\tau_w$ 与 $C_f$。"
\textbf{平板上只需背一个式子：}$\tau_w=\rho U_\infty^2\dfrac{\dd\delta_2}{\dd x}$，其余推导都可现推。
\end{banfa}
\end{kaodian}

\begin{kaodian}{3}{位移厚度 $\delta_1$、动量损失厚度 $\delta_2$ 与形状因子 $H$}
\begin{gongshi}{三种厚度参数的定义（8.5，必背）}
\[ \delta_1=\int_0^{\delta}\left(1-\frac{u}{U}\right)\dd y\ ,\qquad\qquad
   \delta_2=\int_0^{\delta}\frac{u}{U}\left(1-\frac{u}{U}\right)\dd y \]
\[ H=\frac{\delta_1}{\delta_2} \]
\textbf{含义（这是概念题的考点，要会说"物理意义"）：}
\begin{itemize}
  \item $\delta_1$ \textbf{位移厚度}：边界层内速度亏损使通过断面的\emph{质量流量}比理想流体少，
        其数量相当于\emph{壁面外移了 $\delta_1$} 才能恢复同样的流量，
        故又称\emph{排挤厚度}；物理上表示"固体壁面被边界层\"增厚\"了多少"，
        外流相当于绕一个比真实物体\emph{胖 $\delta_1$} 的物体流动。
  \item $\delta_2$ \textbf{动量损失厚度}：边界层内速度亏损使通过断面的\emph{动量通量}减少，
        其数量相当于损失了厚度为 $\delta_2$、速度为 $U$ 的一层流体的动量；
        它\textbf{直接等于摩擦阻力造成的动量损失率}，是卡门方程里的主角。
  \item $H$ \textbf{形状因子}：反映速度剖面的"饱满程度"，
        层流平板 $H\approx2.59$，湍流平板 $H\approx1.3\sim1.4$；
        $H$ \emph{越大剖面越瘦（近壁速度越小）}，也\textbf{越接近分离}，
        故 $H$ 常被用作\emph{分离判据}的经验指标。
\end{itemize}
\textbf{适用条件：}二维不可压缩边界层、不分离；$\delta$ 取同一约定（$0.99U$）。
\end{gongshi}

\textbf{三者关系速查：}$\delta>\delta_1>\delta_2$；层流平板取布拉修斯解时
$\delta_1=1.721\dfrac{x}{\sqrt{\Rey_x}}$，$\delta_2=0.664\dfrac{x}{\sqrt{\Rey_x}}$，$H=2.59$。

\textbf{常考题型：}填空（"位移厚度的物理意义是\underline{\ 因速度亏损使流量减少、相当于壁面外移的厚度\ }"）；
选择（"层流平板边界层的形状因子约为"$\to$2.6）；简答（"比较 $\delta$、$\delta_1$、$\delta_2$ 的定义与物理意义"）。
\textbf{重要度 \xstarfull{3}}\quad\yiju{E1 slide16 填空考"名词定义"、简答考"定义/分类/特点"；E2"掌握基本概念"；E6 教材 8.5 节}\quad\nandu{中}

\begin{banfa}
\textbf{三厚度"一句话区分法"（答题直接用）：}
$\delta$\ ——\emph{几何}厚度：速度恢复到 $0.99U$ 的地方；
$\delta_1$\ ——\emph{流量}亏损：相当于壁面外移多少（管流里算流量修正、势流里修正物面）；
$\delta_2$\ ——\emph{动量}亏损：等于摩擦阻力的动量代价（阻力计算用它）。
\textbf{记忆钩子：}1 管流量、2 管动量、$H$ 管分离。
\end{banfa}
\end{kaodian}

\begin{kaodian}{2}{用假定速度剖面近似求解的步骤}
\textbf{套路（这是"近似解法"的标准流程，务必按顺序会写）：}

\textbf{① 选剖面}：假定 $u/U$ 是 $\eta=y/\delta$ 的简单函数——常用两种（都以边界条件定系数）：
\[ \text{三次多项式：}\ \frac{u}{U}=\frac{3}{2}\eta-\frac{1}{2}\eta^3\ ,\qquad\qquad
   \text{正弦式：}\ \frac{u}{U}=\sin\left(\frac{\pi}{2}\eta\right) \]
三次多项式是用四个条件定出来的：$\eta=0$ 处 $u=0$；
$\eta=1$ 处 $u=U$ 且 $\dfrac{1}{U}\dfrac{\partial u}{\partial \eta}\big|_{\eta=1}=0$；
$\eta=1$ 处 $\dfrac{\partial^2 u}{\partial \eta^2}\big|_{\eta=1}=0$（外缘光滑衔接）。

\textbf{② 算两个厚度}：把剖面代入定义式积分
\[ \delta_1=\delta\!\int_0^1\!\left(1-\frac{u}{U}\right)\dd\eta\ ,\qquad
   \delta_2=\delta\!\int_0^1\!\frac{u}{U}\left(1-\frac{u}{U}\right)\dd\eta \]
三次多项式得 $\delta_1=\dfrac{3}{8}\delta$，$\delta_2=\dfrac{39}{280}\delta$，$H=2.69$。
（注意：各剖面算出的 $H$ 略有差别，但都在 $2.5\sim2.7$，与布拉修斯的 $2.59$ 接近。）

\textbf{③ 写 $\tau_w$}：由牛顿内摩擦定律，壁面切应力只用剖面的\emph{壁面斜率}：
\[ \tau_w=\mu\left.\frac{\partial u}{\partial y}\right|_{y=0}=\mu\frac{U}{\delta}\cdot
   \left.\frac{\dd (u/U)}{\dd \eta}\right|_{\eta=0} \]
三次多项式的壁面斜率为 $1.5$，故 $\tau_w=1.5\mu U/\delta$。

\textbf{④ 代入卡门方程求解}：平板情形用 $\tau_w=\rho U^2\dfrac{\dd \delta_2}{\dd x}$，
把 $\delta_2=C\delta$（$C$ 为常数）与 $\tau_w=A\mu U/\delta$ 代入，得
\[ \frac{A\mu U}{\delta}=\rho U^2 C\frac{\dd\delta}{\dd x}
   \ \Longrightarrow\ \delta\frac{\dd\delta}{\dd x}=\frac{A\nu}{CU}
   \ \Longrightarrow\ \delta=2\sqrt{\frac{A}{C}}\frac{x}{\sqrt{\Rey_x}} \]
三次多项式（$A=1.5$、$C=39/280$）得 $\delta=4.64\dfrac{x}{\sqrt{\Rey_x}}$，$C_f=0.646/\sqrt{\Rey_x}$，
与布拉修斯的精确解（$5.0$、$0.664$）只差约 $7\%$——\textbf{足以说明积分法的工程精度}。

\textbf{⑤ 求阻力}：$C_D=\dfrac{1}{l}\displaystyle\int_0^l C_f\dd x$，或由 $C_D=2\dfrac{\delta_2(l)}{l}$ 直接算出。

\textbf{教材口径（8.6.1 平板层流边界层，书 p168--169，教材原话）：}教材算平板层流边界层用的正是这类"假定剖面＋动量积分"的近似解法，
但假定的剖面\textbf{不是}三次多项式或正弦式，而是\textbf{沿用管流层流（抛物线）剖面}：
"这里假定层流边界层内流速分布和管流中的层流速度分布相同"，即（教材式 8.43）（教材式 8.44）
\[ \frac{u_x}{U_0}=2\frac{y}{\delta}-\left(\frac{y}{\delta}\right)^{2} \]
由牛顿内摩擦定律得（教材式 8.45）$\tau_0=\mu\left.\dfrac{\dd u}{\dd y}\right|_{y=0}=\dfrac{2\mu U_0}{\delta}$；
代入（教材式 8.42）平板边界层基本方程积分得（教材式 8.46）$\delta=\dfrac{5.477x}{\sqrt{\Rey_x}}$，
进而得（教材式 8.47）$\tau_0=0.365\sqrt{\mu\rho U_0^3/x}$、（教材式 8.48）$D=0.73b\sqrt{\mu\rho U_0^3L}$、
（教材式 8.49）$D=C_f\cdot\frac12\rho U_0^2A$、（教材式 8.50）$C_f=1.46/\sqrt{\Rey_L}$（书 p169）。
\textbf{三套结果对照（答题按题目口径选，不要混）：}
\begin{center}
\begin{tabularx}{\textwidth}{@{}lXX@{}}
\toprule
口径 & 假定剖面 & 结果 \\
\midrule
布拉修斯精确解（教材式 8.32、8.35） & 相似性解（数值积分） & $\delta=5.0x/\sqrt{\Rey_x}$，$C_f=0.664/\sqrt{\Rey_x}$ \\
教材动量积分近似（教材式 8.43$\sim$8.50） & 抛物线 $2\eta-\eta^2$ & $\delta=5.477x/\sqrt{\Rey_x}$，$C_f=1.46/\sqrt{\Rey_L}$ \\
本讲义三次多项式 & $\frac32\eta-\frac12\eta^3$ & $\delta=4.64x/\sqrt{\Rey_x}$，$C_f=0.646/\sqrt{\Rey_x}$ \\
\bottomrule
\end{tabularx}
\end{center}
\textbf{教材例题：}（例 8.2，书 p171）用的正是（教材式 8.46）（教材式 8.50）这一近似口径，而非布拉修斯精确解。
另注意（教材式 8.52）（教材式 8.54）是\textbf{湍流边界层}的切应力补充关系式，不适用于本题的层流假定。

\textbf{重要度 \xstarfull{2}}\quad\yiju{E2"第六到第八章只有少量填空/判断/选择，掌握基本概念＋公式套用即可"；E6 教材 8.5 节}\quad\nandu{难}

\begin{banfa}
\textbf{这套步骤的价值在于"不会也能得分的结构分"：}
只要答出"假定剖面 $\to$ 算 $\delta_1$、$\delta_2$ $\to$ 求 $\tau_w$ $\to$ 代入卡门方程 $\to$ 积分得阻力"
这五步，即使积分算错，简答题的分也基本拿到。
\textbf{如果只让估算：}可用正弦式（$\delta_1=0.363\delta$、$\delta_2=0.137\delta$），
它的结果与三次多项式几乎一致。
\end{banfa}
\end{kaodian}

\section{平板边界层计算}

\begin{kaodian}{2}{湍流边界层的速度剖面与公式}
\textbf{是什么：}当 $\Rey_x>5\times10^5$ 后，平板边界层转捩为湍流。
湍流边界层内靠\emph{脉动掺混}传递动量，速度剖面比层流\textbf{饱满}得多，
工程上常用简便的经验公式描述。

\begin{gongshi}{平板湍流边界层的常用公式（8.6）}
\[ \text{速度剖面：}\ \frac{u}{U_\infty}=\left(\frac{y}{\delta}\right)^{1/7}
   \qquad(\text{1/7 次方分布律}) \]
\[ \delta=\frac{0.37x}{\Rey_x^{1/5}}\ ,\qquad\qquad
   C_f=\frac{\tau_w}{\frac{1}{2}\rho U_\infty^2}=\frac{0.0592}{\Rey_x^{1/5}} \]
\[ \text{整板平均：}\ C_D=\frac{0.074}{\Rey_l^{1/5}}\ ,\qquad
   F_d=\frac{1}{2}\rho U_\infty^2\,bl\,C_D \]
\textbf{含义：}用 $\delta$ 作特征尺度、$\Rey_x$ 作无量纲量的幂律拟合法；
$\delta\propto x^{4/5}$（比层流的 $x^{1/2}$ 长得快），$C_f\propto \Rey_x^{-1/5}$（比层流衰减得慢，故摩擦力更大）。
$C_D=0.074/\Rey_l^{1/5}$ 是在 $5\times10^5<\Rey_l<10^7$ 范围内由实验拟合的。\\
\textbf{教材式号：}1/7 次方律即（教材式 8.51）；$\tau_0=0.0225\rho U_0^2\left(\nu/(U_0\delta)\right)^{1/4}$ 即（教材式 8.52）；
$\delta=0.37x/\Rey_x^{1/5}$ 即（教材式 8.53）；$\tau_0=0.0296\rho U_0^2/\Rey_x^{1/5}$ 即（教材式 8.54）；
$D=0.036\rho U_0^2bL/\Rey_L^{1/5}$ 即（教材式 8.55）；$C=0.072/\Rey_L^{1/5}$ 即（教材式 8.56）（书 p169--170）。
\textbf{局部系数换算（教材未直接给局部式，本讲义按教材式推出）：}由（教材式 8.54）$\tau_0=0.0296\rho U_0^2/\Rey_x^{1/5}$，
按 $C_f=2\tau_0/(\rho U_0^2)$ 得\textbf{局部} $C_f=0.0592/\Rey_x^{1/5}$；
由（教材式 8.55）（教材式 8.56）得\textbf{整板} $C_D=0.072/\Rey_L^{1/5}$；二者满足 $C_D=\frac54C_f(L)$。
\textbf{教材原话（关于 0.072 与 0.074，书 p170）：}"与层流边界层比较，当 $\Rey$ 增加时，湍流的 $C$ 要比层流的 $C$ 减小得慢些。
实验表明，上式中的 $0.072$ 改为 $0.074$，则与实验的结果符合得更好。"
即\textbf{教材自身也认可 0.074 更贴合实验}，故本讲义取 $C_D=0.074/\Rey_l^{1/5}$ 与教材（教材式 8.56）相印证，
两者差 $2.7\%$，属取位差异；考试若题目指定教材口径则用 0.072。\\
\textbf{适用条件：}\emph{光滑平板}、\emph{自前缘即全湍流}（或 $\Rey_l$ 很大使层流段可忽略）、
不可压缩、零压梯度、$\Rey_l$ 在 $5\times10^5\sim10^7$；
注意不同教材的系数略有差别（如 $0.074$ 与 $0.072$），\textbf{必须与所用教材配套}；
1/7 次方律只适用于 $\Rey_x<10^7$ 左右，更高雷诺数要用对数律。
\end{gongshi}

\textbf{常考题型：}选择（"湍流边界层厚度沿程按（$x^{ 1/2}$／$x^{4/5}$）增长"$\to x^{4/5}$）；
填空（"湍流边界层常用速度剖面为\underline{\ 1/7 次方分布律\ }"）；
计算（给定条件求 $\delta$、$\tau_w$ 或 $F_d$）。
\textbf{重要度 \xstarfull{2}}\quad\yiju{E2"掌握基本概念＋公式套用即可"；E4 2015 年试题计算题含"平板摩擦阻力比"；E6 教材 8.6 节}\quad\nandu{中}

\begin{banfa}
\textbf{层流 vs 湍流"五对数"对照表（考场速查）：}
\begin{center}
\begin{tabularx}{\textwidth}{@{}lXX@{}}
\toprule
项目 & 层流边界层 & 湍流边界层 \\
\midrule
$\delta/x$ & $5.0/\Rey_x^{1/2}$ & $0.37/\Rey_x^{1/5}$ \\
局部 $C_f$ & $0.664/\Rey_x^{1/2}$ & $0.0592/\Rey_x^{1/5}$ \\
整板 $C_D$ & $1.328/\Rey_l^{1/2}$ & $0.074/\Rey_l^{1/5}$ \\
速度剖面 & 布拉修斯（抛物型，瘦） & 1/7 次方律（饱满） \\
厚度增长 & $\propto x^{1/2}$（缓） & $\propto x^{4/5}$（快） \\
\bottomrule
\end{tabularx}
\end{center}
\textbf{记系数的方法：}层流只有两个数 \textbf{0.664 / 1.328}（后一个是前一个的 2 倍）；
湍流只有两个数 \textbf{0.0592 / 0.074}。
判层/湍只看 $\Rey_l$ 与 $5\times10^5$ 的大小。
\end{banfa}
\end{kaodian}

\begin{kaodian}{2}{混合边界层（前段层流＋后段湍流）的阻力计算}
\textbf{是什么：}实际平板（如机翼、船体）的边界层总是\textbf{前缘一段层流、转捩点之后为湍流}，
这种由层流段与湍流段拼接成的边界层称为\textbf{混合边界层}。
计算思路是"\textbf{分段算、再相加}"，并\emph{扣除}转捩前若已按湍流算所多出的那部分。

\begin{gongshi}{混合边界层的平均阻力系数（8.6）}
\[ C_D=\frac{0.074}{\Rey_l^{1/5}}-\frac{A}{\Rey_l}\quad\text{或写作}\quad
   C_D=\frac{0.074}{\Rey_l^{1/5}}-\frac{1700}{\Rey_l} \]
（后一种为转捩点取 $\Rey_{x,c}=5\times10^5$ 时的常用形式）
\[ \text{整板阻力：}\ F_d=\frac{1}{2}\rho U_\infty^2\,bl\,C_D \]
\textbf{含义：}第一项是把整板都当作湍流算得的值（偏大）；
第二项是\emph{修正项}，扣掉"转捩前本应是层流、却按湍流多算的那部分阻力"，
故 $C_D$ 比"全湍流"小。物理上说明\textbf{前缘保持层流可以减小平板的总摩擦阻力}。\\
\textbf{教材式号：}（教材式 8.57）混合边界层摩擦阻力＝转捩点 $x_c$ 前层流段＋转捩点后湍流段两部分；
（教材式 8.58）其展开式；\textbf{（教材式 8.59a）（教材式 8.59b）}即 $C_{fm}=0.074/\Rey_L^{1/5}-A/\Rey_L$（书 p170--171）。
\textbf{教材假设（书 p170，教材原话）：}"由于混合边界层内流动情况很复杂，在进行计算时，作了两个假设：
一是层流边界层转变为湍流边界层是在 $x_c$ 处\textbf{突然发生}的，\textbf{没有过渡段}；
二是混合边界层的湍流边界层可以看作是从平板首端开始的湍流边界层的一部分。"
\textbf{教材表：}\textbf{表 8.2} 给出修正项 $A$ 的取值（书 p171）：
$\Rey_{x,cr}=10^5\to320$；$3\times10^5\to1050$；$5\times10^5\to1700$；$10^6\to3300$；$3\times10^6\to8700$。
故本讲义"$\Rey_{x,c}=5\times10^5$ 时 $A\approx1700$"与教材表 8.2 \textbf{完全一致}。
\textbf{教材例题：}（例 8.3，书 p171--172）2 m$\times$8 m 平板、$U_0=3$ m/s、$\nu=10^{-6}$ m$^2$/s，
取 $\Rey_{x,cr}=5\times10^5$ 得转捩点 $x_c=0.17$ m，比较"长边顺流"与"短边顺流"两种放法，得摩擦阻力比 $F_1/F_2=0.908$。\\
\textbf{适用条件：}平板、零压梯度、转捩点位置已知（常取 $\Rey_{x,c}=5\times10^5$，此时 $A\approx1700$）；
修正项系数 $A$ 与转捩点位置有关，转捩越靠后、层流段越长，$C_D$ 越小；
若题目直接给出转捩点 $x_c$，也可分段积分
$F_d=\frac12\rho U^2b\!\left(\int_0^{x_c}\!C_{f,\text{层}}\dd x+\int_{x_c}^{l}\!C_{f,\text{湍}}\dd x\right)$。
\end{gongshi}

\textbf{常考题型：}计算（求混合边界层的总阻力，或比较"全层流／全湍流／混合"三种假设的阻力大小）；
选择（"转捩点后移则总摩擦阻力如何变化"$\to$减小）。
\textbf{重要度 \xstarfull{2}}\quad\yiju{E4 2015 年试题计算题含"平板摩擦阻力比"；E2"掌握基本概念＋公式套用即可"；E6 教材 8.6 节}\quad\nandu{中}

\begin{banfa}
\textbf{"三算一比较"套路：}
① 先算 $\Rey_l$ 判断是否可能全层流（$\Rey_l<5\times10^5$ 时整个板都是层流，\emph{不要}用混合公式）；
② 若 $\Rey_l>5\times10^5$，判题意："全湍流"用 $0.074/\Rey_l^{1/5}$，"混合"再减 $A/\Rey_l$；
③ 代 $F_d=\frac12\rho U^2blC_D$（\emph{单面}，双面乘 2）；
④ \textbf{比大小结论要背}：同 $\Rey_l$ 下 $C_{D,\text{全层}}>C_{D,\text{全湍}}$ 还是相反？
—— 由 $\Rey_l^{-1/2}$ 与 $\Rey_l^{-1/5}$ 比较：$\Rey_l>10^5$ 时 $\Rey_l^{-1/2}<\Rey_l^{-1/5}$，
\textbf{所以 $C_{D,\text{层}}<C_{D,\text{湍}}$，即层流摩擦阻力更小}（这也正是"层流翼型"存在的理由）。
\end{banfa}

\begin{yicuo}
\textbf{常见陷阱：}湍流边界层"速度剖面饱满、近壁速度梯度大"，
所以\textbf{湍流的摩擦阻力大于层流}（$C_f$ 大）；但湍流\textbf{更不易分离}，\emph{形状阻力小}。
两句话必须一起说：\textbf{湍流＝摩擦阻力大＋抗分离能力强}。
只答前半句或只答后半句都会丢分。
\end{yicuo}
\end{kaodian}

\section{边界层的分离及减阻}

\begin{kaodian}{3}{边界层分离的物理机制与分离点判据}
\textbf{是什么：}边界层在\textbf{逆压梯度}区段内，近壁\emph{动能不足}的流体被压强顶住，
速度先减小到零，继而出现\textbf{倒流}，使边界层脱离壁面的现象称为\textbf{分离}，
发生分离的位置叫\textbf{分离点}。

\textbf{物理机制（简答题的标准答法，要按因果链写）：}
\begin{enumerate}
  \item 逆压梯度 $\dfrac{\dd p}{\dd x}>0$（等价 $\dd U/\dd x<0$）：流体沿流向要\emph{顶着压强爬坡}，
        动能不断转化为压力能，速度下降；
  \item 边界层内近壁流体本来\emph{动能最小}（壁面 $u=0$；靠黏性从外部缓慢取得动量），
        先于外层被"刹停"；
  \item 壁面处速度梯度 $\left.\dfrac{\partial u}{\partial y}\right|_{y=0}$ 由正变零、再由零变负，
        近壁出现\textbf{回流}；
  \item 回流把边界层"顶"离壁面，形成\textbf{分离区（尾流区）}与旋涡，
        主流绕分离区流过，形成\emph{尾涡}。
\end{enumerate}

\textbf{教材原话（书 p172--173，8.7.1 边界层分离）：}"流速为零，压强最大的点，称为\textbf{停滞点}或\textbf{驻点}。"
／"到贴近圆柱面的 C 点，流速将为零。流体质点在 C 点停滞下来，形成新的驻点。由于流体不可压缩，
继续流来的流体质点被迫脱离原来的流线，沿着另一条流线流去，如图中的 CE 线，从而使边界层脱离了圆柱面，
这种现象即为\textbf{边界层的分离现象}，C 点称为\textbf{分离点}。"
／"边界层分离后，在边界层与圆柱面之间，由于分离点下游的压强大，从而使流体发生反向回流，形成旋涡区。
在绕流物体边界层分离点下游形成的旋涡区称为\textbf{尾流}。分离点的位置是不固定的，它和流体所绕物体的形状、
粗糙程度、流动的雷诺数等有关。如流体遇到固体表面的锐缘时，分离点就在锐缘处。另外，边界层的分离还与来流和物体的相对方向有关。"
／"当平板与来流方向平行放置时，边界层不会发生分离。但当平板与来流方向垂直放置时，则必然在平板两端产生分离。"
\textbf{教材提示：}教材 8.7.1 节\textbf{没有}给分离点判据的式号，该节以文字叙述为主；
本讲义公式盒题名中的"8.7"表示\emph{节}号，\emph{不表示}教材式号，特此说明以免误引。

\begin{gongshi}{分离点判据（8.7 节，必背）}
\[ \text{分离点：}\quad\left.\frac{\partial u}{\partial y}\right|_{y=0}=0 \]
\[ \text{分离前：}\ \left.\frac{\partial u}{\partial y}\right|_{y=0}>0\ ;\qquad
   \text{分离后：}\ \left.\frac{\partial u}{\partial y}\right|_{y=0}<0 \]
\[ \text{壁面切应力：}\ \tau_w=\mu\left.\frac{\partial u}{\partial y}\right|_{y=0}
   \ \Longrightarrow\ \text{分离点处}\ \tau_w=0 \]
\textbf{含义：}分离点即壁面切应力为零的点（速度剖面在壁面处的切线变竖直）。
该判据把"分离"这件几何上难观测的事，转化为\textbf{壁面切应力过零}这一可计算的条件。\\
\textbf{适用条件：}二维、不可压缩、连续介质；流动必须先出现\textbf{逆压梯度}；
层流边界层可用该判据直接判断；湍流边界层由于脉动掺混带来的动量补充，
需要用"壁面定律"或半经验判据（如形状因子 $H$ 的经验阈值）修正。
\end{gongshi}

\begin{jielun}
\textbf{分离的两个必要条件缺一不可：}\textbf{逆压梯度（$\dd p/\dd x>0$）＋ 黏性}。
—— 理想流体无黏，壁面处\emph{允许滑移}、$\partial u/\partial y$ 不趋于无穷大，
且没有近壁低速层，\textbf{无论如何都不会分离}；
—— 有黏但没有逆压梯度（如平板），速度剖面沿程相似放大，\textbf{也不会分离}。
因此"边界层分离是黏性作用与逆压梯度共同作用的结果"是标准结论。
\end{jielun}

\textbf{常考题型：}简答（"试述边界层分离的物理原因与分离条件"——\textbf{高频}）；
填空（"边界层分离点的条件是\underline{\ $\partial u/\partial y|_{y=0}=0$\ }"）；
判断（"理想流体也会发生边界层分离"$\to$错；"只有逆压梯度就会分离"$\to$错）。
\textbf{重要度 \xstarfull{3}}\quad\yiju{E1 slide16 简答考"定义、分类、特点"式问题；E2"掌握基本概念"；E6 教材 8.7 节}\quad\nandu{中}

\begin{banfa}
\textbf{简答题"四句话"模板（背下来直接默写）：}
① "在逆压梯度区，压强沿程升高，边界层内流体动能不断消耗、速度减小"；
② "近壁处流体质点动能最小，先减速到零"；
③ "此后出现回流，$\left.\partial u/\partial y\right|_{y=0}=0$ 处即为分离点"；
④ "分离形成尾涡区，使物体前后压差增大，产生形状阻力（压差阻力）。"
\textbf{顺手补一句：}"理想流体无黏、无边界层，故不会分离——分离是黏性现象。"
\end{banfa}

\begin{yicuo}
\textbf{易错 4：}分离的根源是\textbf{逆压梯度＋黏性，两者缺一不可}。
只答"有逆压梯度就分离"是错的（平板有黏性、无逆压，也不分离）；
只答"黏性导致分离"也不完整（无逆压梯度同样不分离）。
理想流体无黏，\emph{永远不分离}——这是判断题最爱设的坑。
\end{yicuo}
\end{kaodian}

\begin{kaodian}{3}{分离与形状阻力、尾涡的关系；摩擦阻力与形状阻力的区别}
\textbf{两种阻力（对比记忆，选择/填空常考）：}
\begin{center}
\begin{tabularx}{\textwidth}{@{}lXX@{}}
\toprule
项目 & \textbf{摩擦阻力} & \textbf{形状阻力（压差阻力）} \\
\midrule
产生原因 & 流体的黏性在壁面上产生的切应力 & 分离后尾涡区使物体前后压强不平衡 \\
正比于 & \textbf{湿表面积}（与表面接触的面积） & \textbf{迎流面积}与尾流区尺寸 \\
与分离关系 & 分离前后都存在 & \textbf{只有分离后才显著} \\
典型例 & 平板平行来流（几乎只有摩擦阻力） & 平板垂直来流、圆柱、钝体（几乎只有形状阻力） \\
减小方法 & 减小湿面积、层流化、微沟槽、降黏 & 流线型化、延缓分离 \\
\bottomrule
\end{tabularx}
\end{center}

\textbf{总阻力：}$F_D=F_{\text{摩擦}}+F_{\text{形状}}$，两者之比随物体形状变化很大；
对\emph{流线型}物体摩擦阻力占主导，对\emph{钝体}形状阻力占主导。

\textbf{尾涡与能量损失的物理图像：}分离后尾流区是低压的旋涡区，
旋涡不断被主流带走、又不断生成，把主流的一部分机械能不可逆地转化为热能，
既恶化了物体前后压差（阻力增大），又引起机械振动与噪声。
工程上"避免在重要部位出现分离"就是为此。

\textbf{常考题型：}选择（"下列哪种阻力与湿表面积成正比"$\to$摩擦阻力）；
填空（"物体的总阻力由\underline{\ 摩擦阻力\ }和\underline{\ 形状阻力\ }组成"）；
判断（"分离后摩擦阻力消失"$\to$错）。
\textbf{重要度 \xstarfull{3}}\quad\yiju{E1 slide16 选择/判断考"概念辨析"、填空考"名词定义"；E2"掌握基本概念"；E6 教材 8.7 节}\quad\nandu{易}

\begin{banfa}
\textbf{一句口诀定"谁能减小谁"：}"摩擦靠面积，形状靠尾涡"——
问"微沟槽减阻减的是什么"$\to$摩擦阻力（湿面积不变但切应力降）；
问"流线型化减的是什么"$\to$形状阻力（延缓分离、缩小尾涡）。
两者不要互串。
\end{banfa}
\end{kaodian}

\begin{kaodian}{2}{绕流阻力的分解与圆球自由沉降（教材 8.7.2 新增）}
\textbf{为什么补这一节：}教材 8.7.2 把"绕流阻力"单列一节并给出可直接套用的公式与\textbf{例 8.4}，
本讲义原稿只有摩擦阻力／形状阻力的\emph{概念对比}，缺\textbf{公式}与\textbf{颗粒沉降}这一支。
历年真题中"颗粒（钢球）沉降"类计算反复出现，故按教材补出。

\begin{gongshi}{绕流阻力的分解与合成（8.7.2，必背）}
\[ D=D_f+D_p \]
\[ D_f=C_f\frac{1}{2}\rho U_\infty^2A_s\ ,\qquad
   D_p=C_p\frac{1}{2}\rho U_\infty^2A_p\ ,\qquad
   D=C_D\frac{1}{2}\rho U_\infty^2A \]
\textbf{教材式号：}（教材式 8.60）$D_f$；（教材式 8.61）$D_p$；（教材式 8.62）$D$（书 p173）。
其中 $A_s$ 为\textbf{切应力作用的面积}；$A_p$ 为物体与流速方向垂直的\textbf{迎流投影面积}；$A=A_p$。
\textbf{教材原话：}"绕流阻力 $D$ 包括摩擦阻力 $D_f$ 和压差阻力 $D_p$ 两部分"；
"压差阻力主要取决于物体的形状，所以又称为\textbf{形状阻力}"；
$C_D$"主要取决于雷诺数，并和物体的形状、表面粗糙度以及来流的紊动强度有关……尚无法由理论计算得出，多由实验确定"；
教材图 8.13 给出\textbf{圆球、圆盘及无限长圆柱}的阻力系数实验曲线。\\
\textbf{适用条件：}不可压缩定常绕流；$A_s$ 与 $A_p$ 的取法必须与所用 $C_f$、$C_p$ 的定义配套；
$C_D$ 是 $\Rey$ 的函数（圆球用相对来流速度 $U_\infty$ 与球径）。
\end{gongshi}

\begin{gongshi}{圆球自由沉降速度（8.7.2，必背）}
\[ u_f=\sqrt{\frac{4(\rho_s-\rho)gd}{3C_D\rho}} \]
\textbf{教材式号：}（教材式 8.63）（书 p174）。
\textbf{含义（教材原话）：}"圆球的重量与所受的浮力和阻力达到平衡，圆球作等速沉降，
其速度称为\textbf{自由沉降速度}，用 $u_f$ 表示。"平衡式为 $G=B+D$，
其中 $G=\frac16\pi d^3\rho_s g$、$B=\frac16\pi d^3\rho g$、$D=C_D\frac12\rho u_f^2\cdot\frac{\pi d^2}{4}$；
$\rho_s$ 为球体密度、$\rho$ 为流体密度、$d$ 为球径。\\
\textbf{$C_D$ 的分段经验式（教材口径，书 p175）：}层流段 $C_D=24/\Rey$（教材该段首行适用范围数字 OCR 无法判读）；
过渡段 $C_D=13/\sqrt{\Rey}$，"圆球呈摆动状态下沉，绕流属于过渡状态"；
湍流段 $C_D\approx0.45$，"当 $\Rey=10^3\sim2\times10^5$ 时，圆球脱离铅垂线，盘旋下沉，绕流属于湍流状态"。
\textbf{必须试算（教材原话）：}"计算自由沉降速度，因为 $u_f$ 与 $\Rey$ 有关，而 $\Rey$ 中又包含待求值 $u_f$，
所以一般要经过多次试算才能求得……可以先假定 $\Rey$ 的范围，然后再验算 $\Rey$ 是否与假定的一致"。
\textbf{悬浮速度（教材原话）：}"当 $u=u_f$ 时，$u_s=0$，则圆球悬浮在流体中，呈悬浮状态，
这时流体上升的速度 $u$ 称为圆球的\textbf{悬浮速度}，它的数值与 $u_f$ 相等，但意义不同。
自由沉降速度是圆球自由下降时所能达到的最大速度，而悬浮速度是流体上升速度能使圆球悬浮所需的最小速度。"\\
\textbf{教材例题：}（例 8.4，书 p175）炉膛烟气上升速度 $u=0.5$ m/s、$\rho=0.2$ kg/m$^3$、$\nu=230\times10^{-6}$ m$^2$/s，
煤粉 $d=0.1$ mm、$\rho_s=1300$ kg/m$^3$，试算得 $C_D=110.6$、$u_f=0.278$ m/s；
教材结论原话："因为 $u=0.5$ m/s $>u_f=0.278$ m/s，所以煤粉颗粒将被烟气流带走，不会沉降。"
（历年真题"钢球沉降测黏度／颗粒沉降"一类题即此型，判据是 $u_f$ 与来流速度 $u$ 的比较。）
\end{gongshi}

\textbf{常考题型：}填空（"绕流阻力由\underline{\ 摩擦阻力\ }和\underline{\ 压差阻力\ }两部分组成"）；
选择（"圆球自由沉降速度与 $\sqrt{d}$、$\sqrt{\rho_s-\rho}$ 成正比"）；
计算（给 $d$、$\rho_s$、$\rho$、$\nu$ 求 $u_f$ 或判断能否沉降——\textbf{必须试算}）。
\textbf{重要度 \xstarfull{2}}\quad\yiju{E2"第六到第八章只有少量填空/判断/选择，掌握基本概念＋公式套用即可"；E6 教材 8.7.2 节}\quad\nandu{中}

\begin{banfa}
\textbf{沉降题三步走：}① 假定 $\Rey$ 区间（先猜层流，取 $C_D=24/\Rey$）$\to$
② 代 $u_f=\sqrt{4(\rho_s-\rho)gd/(3C_D\rho)}$ 求 $u_f$ $\to$ ③ 用求得的 $u_f$ 反算 $\Rey$，验算是否落在假定区间，不符则换段重算。\\
\textbf{结论句一律写成"比较 $u$ 与 $u_f$"：}$u>u_f$ 被带走、$u<u_f$ 会沉降、$u=u_f$ 悬浮。
\textbf{打通概念题与公式题：}"摩擦阻力 $\propto$ 湿面积（切应力作用面积 $A_s$）""压差阻力 $\propto$ 迎流投影面积 $A_p$"——
答题时把教材符号 $A_s$、$A_p$ 一并写出，比只写汉字更稳。
\end{banfa}
\end{kaodian}

\begin{kaodian}{3}{减阻措施（简答题高频）}
\textbf{总纲：}减阻要么\textbf{减小摩擦阻力}（管好边界层内的切应力与湿面积），
要么\textbf{减小形状阻力}（管住分离与尾涡）。措施按这个二分法组织，答题就不会漏。

\begin{center}
\begin{tabularx}{\textwidth}{@{}lXX@{}}
\toprule
措施 & 做法 & 减的是什么／原理 \\
\midrule
① 流线型化、优化型线 & 采用水滴形、细长比合适的型线，使逆压梯度平缓 &
减形状阻力：推迟分离、缩小尾涡；同时尽量减小湿面积 \\
② 减小逆压梯度 & 设计时使 $\dd p/\dd x$ 的逆压段尽量短、峰值尽量小（如尾部缓收） &
减形状阻力：给边界层留更多动能，推迟分离 \\
③ 边界层抽吸（吸气） & 通过多孔壁把近壁低速、已将分离的流体吸入内部 &
同时减两种阻力：把低动能流体移走，剖面重新饱满；是延缓分离最有效的手段之一 \\
④ 使边界层提前转捩为湍流 & 加绊线、粗糙带或小涡流发生器，令转捩点前移 &
减形状阻力：湍流掺混强、抗逆压能力好，\textbf{推迟分离}；代价是摩擦阻力略增 \\
⑤ 微沟槽（肋条）表面减阻 & 沿流向开 $\mathrm{\mu m}$ 级的顺流向沟槽（仿鲨鱼皮） &
减摩擦阻力：抑制流向涡与湍流猝发，可减摩擦阻力数个百分点 \\
⑥ 添加高分子聚合物 & 在水中加入少量高分子聚合物（如聚丙烯酰胺） &
减摩擦阻力：抑制湍流脉动与近壁相干结构，降低湍流切应力 \\
⑦ 控制表面粗糙度、柔性壁面 & 抛光、避免不必要的凸起与台阶 &
减摩擦阻力：避免局部粗糙把层流底层"戳破"而转为水力粗糙 \\
⑧ 边界层加速（吹气／缝翼） & 向将分离处注入高速流体；机翼内部设喷气气源，或前缘加设缝翼形成"喷嘴翼" &
同时减两种阻力：给近壁流体补充动能、推迟分离（\textbf{教材专列的三大方法之一，原稿缺失}） \\
\bottomrule
\end{tabularx}
\end{center}

\textbf{教材口径（8.7.3，书 p175--176，简答题按教材答最稳）：}教材把减阻归为\textbf{三条控制边界层的方法}——
\begin{center}
\begin{tabularx}{\textwidth}{@{}lXX@{}}
\toprule
教材序号 & 做法（教材原话要点） & 对应本讲义上表 \\
\midrule
(1) 流线型 & "将被绕流物体的外形设计成流线型……物体后部曲率越大，分离越早，尾流越粗，压差阻力相应越大" &
① \\
(2) 边界层加速 & "向边界层注入高速流体……图 8.15（a）是在机翼内部设置一喷气气源，将高速射流从边界层将要分离处喷入边界层。图 8.15（b）是在机翼前缘处加设一缝翼，它与机翼之间形成一喷嘴翼" &
\textbf{⑧（原稿缺失，本次补出）} \\
(3) 边界层的吸收 & "在边界层易分离处设置一窄缝，在机翼内的抽气装置把欲滞止的空气经该缝抽走，从而防止了分离……这种方法还可以使边界层的层流到湍流的\textbf{转捩点后移}，达到减小摩擦阻力的效果" &
③、⑦ \\
\bottomrule
\end{tabularx}
\end{center}

\textbf{教材两个关键结论（原话，书 p175--176）：}"绕流阻力中的压差阻力和摩擦阻力的主次取决于雷诺数。
对于流体绕流圆柱体，当雷诺数较小时，压差阻力占总阻力的 1/3；当雷诺数增大时，压差阻力占到总阻力的一半；
当 $\Rey=200$ 时，压差阻力增至总阻力的 75\%"（该段末句"当 $\Rey=10\sim10$ 时，总阻力主要是压差阻力"中
的\textbf{指数被 OCR 丢失，无法判读}）；／"一般来说，层流边界层的摩擦阻力比湍流边界层小，是湍流时的 $1/5\sim1/6$。
为了减小摩擦阻力，应采用小的物面粗糙度，使物面上的层流边界层尽可能长。"／"边界层分离形成的压差阻力会使绕流阻力增大，
而升力骤减，导致叶片式流体机械的运行效率下降。"

\begin{tishi}
\textbf{订正（原表述 $\to$ 新表述）：}原稿减阻表把"使边界层提前转捩为湍流"（④）列为一条减阻措施，
\textbf{教材 8.7.3 没有这条}；教材的 (3) 边界层吸收讲的是\textbf{转捩点后移}（把层流段拉长以减小摩擦阻力），
与本讲义④的方向\emph{正好相反}。两者都成立，但\textbf{前提必须说清}：
④（提前转捩）靠湍流掺混抗逆压，\emph{推迟分离、减形状阻力}，代价是摩擦阻力增大；
教材 (3)（抽吸使转捩点后移）\emph{延长层流段、减摩擦阻力}，它与"层流边界层摩擦阻力只有湍流的 $1/5\sim1/6$"这条教材结论配套。
故答题时按"以形状阻力为主用④，以摩擦阻力为主用教材的 (3)"择一展开，切勿把两者的方向混说。
另：教材 (2)"边界层加速"原稿完全没有，属\textbf{漏项}，已补为上表第⑧条。
\end{tishi}

\textbf{权衡（务必写上一句）：}③④⑤⑦ 主要减\textbf{摩擦}阻力，②③④ 主要减\textbf{形状}阻力；
而④（提前转捩）是\textbf{以摩擦阻力增大为代价换取分离推迟}，
因此在\emph{以形状阻力为主}的场合（机翼大迎角、钝体绕流）才划算，
在以摩擦阻力为主的场合（机翼巡航、平板）反而不利——工程上要按工况权衡。

\textbf{常考题型：}\textbf{简答（高频）}："简述减小流体阻力的措施"／"如何防止边界层分离"；
选择（"下列措施中不是减小摩擦阻力的是"）；
判断（"使边界层提前转捩总能减小总阻力"$\to$错，需分场合）。
\textbf{重要度 \xstarfull{3}}\quad\yiju{E1 slide16 简答考"定义、分类、特点"式问题（E1 slide16 明确简答题的考法）；E2"掌握基本概念"；E6 教材 8.7 节}\quad\nandu{易}

\begin{banfa}
\textbf{答案要"成对出现"才拿满分：}每写一条措施，就补一句它\emph{减哪种阻力、代价是什么}。
写成"措施＋减阻类型＋代价"三列，简答题 4 分基本能拿满 3--4 分。
\textbf{最容易被漏掉的两条：}③\textbf{边界层抽吸}（教材上专门讲，务必写上）与
⑥\textbf{添加高分子聚合物}（教材提到的减阻新手段）。
\textbf{提防判断题陷阱：}"提前转捩总能减阻"$\to$\textbf{错}，它只在形状阻力为主的场合有利。
\end{banfa}

\begin{yicuo}
\textbf{易错 5：}摩擦阻力 $\propto$ \textbf{湿表面积}；形状阻力 $\propto$ \textbf{迎流面积与尾流区}。
这两个成正比的对象是反着来的：船体/机翼越"流线"，湿面积越大（摩擦阻力上升）但尾涡越小（形状阻力下降），
所以工程上存在\textbf{最优长细比}——不是越流线越好。
\end{yicuo}
\end{kaodian}

\section{湍流的基本方程}

\begin{kaodian}{3}{雷诺时均法与雷诺方程}
\textbf{思路（湍流为什么不能直接用 N-S 方程）：}湍流是三维、有旋、非定常的随机流动，
瞬时量随时间和空间剧烈脉动；直接求解 N-S 方程需要分辨所有尺度的涡，
在工程雷诺数下计算量大到不可行。工程上改用\textbf{时均化}方法，
只求\emph{平均值}，把脉动的影响归入附加应力。

\textbf{雷诺分解（8.8）：}把瞬时量分解为时均值与脉动量之和
\[ u=\bar u+u'\ ,\qquad v=\bar v+v'\ ,\qquad w=\bar w+w'\ ,\qquad p=\bar p+p' \]
其中 $\bar u=\dfrac{1}{T}\displaystyle\int_0^T u\,\dd t$ 为\textbf{时均速度}，
$T$ 取远大于脉动周期、远小于流动宏观特征时间；
由定义立即得 $\overline{u'}=0$，即\textbf{脉动量的时均值恒为零}。

\textbf{时均化运算规则（填空小考点）：}
\[ \overline{\bar u}=\bar u\ ,\qquad \overline{u'}=0\ ,\qquad
   \overline{\bar u\,v'}=\bar u\,\overline{v'}=0\ ,\qquad
   \overline{u'+v'}=\overline{u'}+\overline{v'}=0 \]
\[ \overline{\bar u\,\bar v}=\bar u\,\bar v\ ,\qquad
   \overline{u'^2}>0\ (\text{脉动量的平方时均值不为零！}) \]

\textbf{教材的时均化运算法则（8.8.1，书 p177，教材原文共 8 条）：}
\begin{enumerate}
  \item $f\pm g$ 的时均值 $=\bar f\pm\bar g$；
  \item $\overline{af}=a\bar f$（$a$ 为常数）；
  \item $\overline{\bar f}=\bar f$（时均量的时均值仍为其本身）；
  \item $\overline{f'}=0$（脉动量的时均值均为零）；
  \item $\overline{\bar f\bar g}=\bar f\bar g$；
  \item $\overline{\bar f g'}=0$；
  \item $\overline{fg}=\bar f\bar g+\overline{f'g'}$（\textbf{填空高频}）；
  \item $\overline{\dfrac{\partial^n f}{\partial s^n}}=\dfrac{\partial^n\bar f}{\partial s^n}$（时均与求导可交换）。
\end{enumerate}
\textbf{教材式号：}湍流的连续性方程为\textbf{（教材式 8.64a）（教材式 8.64b）}（书 p177）。

\begin{gongshi}{雷诺方程（时均 N-S 方程，8.8，必背）}
\[ \text{连续：}\ \frac{\partial \bar u}{\partial x}+\frac{\partial \bar v}{\partial y}+\frac{\partial \bar w}{\partial z}=0 \]
\[ \rho\left(\frac{\dd \bar u}{\dd t}\right)
   =\rho X-\frac{\partial \bar p}{\partial x}
   +\mu\grad^2\bar u
   +\frac{\partial}{\partial x}\!\left(-\rho\overline{u'u'}\right)
   +\frac{\partial}{\partial y}\!\left(-\rho\overline{u'v'}\right)
   +\frac{\partial}{\partial z}\!\left(-\rho\overline{u'w'}\right) \]
\textbf{含义：}把雷诺分解代入 N-S 方程再做时均化，
左边仍保留时均速度的惯性项与压力项、黏性项，
但\textbf{多出三项} $-\rho\overline{u'_iu'_j}$ 的偏导——这就是\textbf{雷诺应力}
（又称\textbf{湍流附加应力}或\textbf{湍流应力}）。
它的物理来源是：脉动使流体微团在断面之间\textbf{交换动量}，
宏观上等效为一组附加应力。\emph{湍流的本质＝脉动导致的动量交换}。\\
\textbf{教材式号：}由 $x$ 向 N-S 方程依次得（教材式 8.65）（教材式 8.66）（教材式 8.67）（教材式 8.68）（教材式 8.69），
最终\textbf{（教材式 8.70）为三个方向上的雷诺方程}（书 p177--178）。
\textbf{教材原话（书 p178）：}"将雷诺方程和 N-S 方程比较，发现雷诺方程中增加了 6 个与脉动相关的应力，
称为\textbf{雷诺应力}或\textbf{湍动应力}。同黏性应力类似，雷诺应力也可用对称矩阵表示。"
\textbf{教材推导起点（书 p176）：}教材先说明"理论上描述瞬时流动的 N-S 方程同样适用于湍流。但由于湍流运动相当复杂，
求解湍流瞬时流动的全部过程，既不必要也不可能……有意义的是湍流过程总体的统计特性。最重要的同时也是最简单的统计特征值是平均值。"\\
\textbf{适用条件：}不可压缩牛顿流体、定常时均（$\partial\bar u/\partial t=0$）；
\emph{时均}方程组只描述平均流动，脉动信息已被归入雷诺应力。
\end{gongshi}

\textbf{不可压缩时均方程组的封闭性问题（必答结论）：}
未知量有 $\bar u,\bar v,\bar w,\bar p$ 共 4 个，但\emph{新增的}雷诺应力有
$\overline{u'u'}$、$\overline{v'v'}$、$\overline{w'w'}$、$\overline{u'v'}$、$\overline{u'w'}$、$\overline{v'w'}$
共 \textbf{6 个独立分量}（由对称性）。
3 个动量方程＋1 个连续方程＝4 个方程，却有 $4+6=10$ 个未知量，
\textbf{方程数少于未知量数，方程组不封闭}。
必须补充\textbf{湍流模型}（把雷诺应力与平均速度场联系起来，见 8.9 节）才能求解。

\textbf{常考题型：}填空（"雷诺方程比 N-S 方程多出的项称为\underline{\ 雷诺应力\ }"；
"脉动量的时平均值等于\underline{\ 零\ }"）；
选择（"雷诺应力的物理来源"$\to$脉动引起的动量交换）；
简答（"写出雷诺方程并说明雷诺应力的含义"）。
\textbf{重要度 \xstarfull{3}}\quad\yiju{E1 slide33 第八章＝"选择填空题考点"；E1 slide16 填空考"名词定义"；E2"掌握基本概念"；E6 教材 8.8 节}\quad\nandu{中}

\begin{banfa}
\textbf{"多三行"记忆法：}雷诺方程＝N-S 方程＋三行 $-\rho\overline{u'_iu'_j}$ 的散度项。
答题时先默写 N-S 方程，再补一句"对瞬时量做时均化后，多出雷诺应力项"即可。
\textbf{封闭性问题的标准三句：}"雷诺应力有 6 个独立分量，方程只有 4 个，
未知量多于方程，故方程组不封闭，必须引入湍流模型——这就是所有湍流模型的出发点。"
\end{banfa}

\begin{yicuo}
\textbf{易错 6：}不要把 N-S 方程说成"只适用于层流"。
\textbf{N-S 方程本身对层流和湍流（瞬时量）都成立}；
湍流之所以要另立方程组，不是因为 N-S 方程失效，
而是因为\emph{直接数值求解瞬时量的代价太大}，故改用\emph{时均化}的雷诺方程。
这是选择题与判断题的经典设坑点，也是本章与第 5 章"两种流态"最容易混起来的地方。
\end{yicuo}
\end{kaodian}

\begin{kaodian}{2}{普朗特混合长度理论与湍流切应力}
\textbf{思路（第一个实用的湍流模型）：}雷诺应力是未知的，普朗特设想用力学类比把它表达成
\emph{平均速度梯度}的函数——湍流脉动像一个"微小流体团"跳出特征距离后再与周围掺混合并。

\begin{gongshi}{混合长度假设与湍流切应力（8.8）}
\[ \tau_t=-\rho\overline{u'v'}=\rho\,l^2\left(\frac{\dd \bar u}{\dd y}\right)^2
   \qquad(\text{取绝对值})\ ,\qquad l=\kappa y \]
\[ \text{总切应力：}\ \tau=\tau_{\text{黏}}+\tau_t
   =\mu\frac{\dd \bar u}{\dd y}+\rho\,l^2\left(\frac{\dd \bar u}{\dd y}\right)^2
   =\rho\left(\nu+l^2\left|\frac{\dd \bar u}{\dd y}\right|\right)\frac{\dd \bar u}{\dd y} \]
\textbf{含义：}$l$ 称\textbf{混合长度}，表示流体微团在丧失原有属性前横向运动的\emph{特征距离}，
可理解为湍流"掺混的尺度"；$\kappa$ 是卡门常数（$\kappa\approx0.4$）。
上式把未知的雷诺应力化归为\emph{已知的平均速度梯度}，
从而\textbf{使时均方程组可以封闭}——这是零方程湍流模型的核心。\\
\textbf{教材式号：}（教材式 8.71）$\tau_y=\rho\overline{u'_x u'_y}\,\dd A_y$；（教材式 8.72）$\tau=-\rho\overline{u'_x u'_y}$；
（教材式 8.73）涡黏性模型（布辛涅斯克 Boussinesq 提出）$-\rho\overline{u'_iu'_j}=\nu_t\left(\partial u_i/\partial x_j+\partial u_j/\partial x_i\right)$；
\textbf{（教材式 8.74a）（教材式 8.74b）}$\tau=\rho l^2\left|\dd\bar u/\dd y\right|\left(\dd\bar u/\dd y\right)$；（教材式 8.75）$l=ky$；
（教材式 8.76）$\tau=\mu\,\dd\bar u/\dd y+\rho l^2\left(\dd\bar u/\dd y\right)^2$；
（教材式 8.77）$u=\dfrac1k\sqrt{\dfrac{\tau_0}{\rho}}\ln y+C$（边界附近湍流流速的\textbf{对数分布}）（书 p179--182）。
\textbf{教材原话：}"上式即湍动切应力与脉动流速之间的关系式，是雷诺在 1895 年提出的"（书 p179）；
卡门常数的教材取值："实验表明 $k=0.36\sim0.435$，常取 $k=0.4$；$y$ 为从壁面算起的横向距离"（书 p181）；
"层流时，流体质点无横向混掺现象，$l=0$，流层间只有黏性切应力，$\tau=\mu\dfrac{\dd u}{\dd y}$"；
"在雷诺数相当大的情况下，（黏性切应力）比湍动切应力小得多，可以忽略不计，流层间的切应力只有湍动切应力"（书 p181）。
\textbf{涡黏性模型与黏度的区别（教材原话，书 p180）：}"湍动黏度 $\nu_t$ 和黏度 $\nu$ 之间有本质的区别：
湍动黏度 $\nu_t$ 反映\textbf{湍动特性}，与流动状况和边界条件密切相关；黏度 $\nu$ 表示的是流体的一种物理属性，
其值取决于流体的性质，而与流动状况无关。"（\textbf{填空/选择高频}）\\
\textbf{教材对混合长度理论的评价（原话，书 p182）：}"混合长度理论使雷诺方程组封闭，且能结合实验解决一些实际问题，
具有重要意义。但由于其属于湍流的半经验理论，有些假设还不是很严格，如，假定流体质点要经过一定距离才发生其他质点相混掺，
这与实际混掺是一连续过程不符。"\\
\textbf{适用条件：}\emph{局部}平衡的剪切湍流（近壁区、自由剪切层）适用；
$\tau_t$ 只取大小（故写成梯度平方或绝对值形式），
\emph{不能反映上游历史效应与逆压梯度的输运影响}；
$l=\kappa y$ 只在贴近壁面处成立，远离壁面需另行给定。
\end{gongshi}

\textbf{黏性切应力与湍流切应力的对比（填空/选择常考）：}
\begin{center}
\begin{tabularx}{\textwidth}{@{}lXX@{}}
\toprule
项目 & 黏性切应力 $\tau_{\text{黏}}$ & 湍流切应力 $\tau_t$ \\
\midrule
表达式 & $\mu\,\dfrac{\dd \bar u}{\dd y}$ & $\rho\,l^2\left(\dfrac{\dd \bar u}{\dd y}\right)^2$ \\
物理机制 & \textbf{分子}热运动引起的动量交换 & \textbf{流体质点（微团）}脉动的动量交换 \\
与梯度关系 & \textbf{线性} & \textbf{二次}（梯度越大增长越快） \\
主导区域 & 黏性底层（紧贴壁面） & 湍流核心区 \\
是否耗散 & 是，黏性耗散 & 主要是输运，最终仍由小尺度涡耗散 \\
\bottomrule
\end{tabularx}
\end{center}

\textbf{常考题型：}填空（"普朗特混合长度理论中 $\tau_t=$ \underline{\ }"）；
选择（"黏性底层中主要起作用的是（黏性／湍流）切应力"$\to$黏性）；
简答（"简述雷诺应力与黏性切应力的区别"）。
\textbf{重要度 \xstarfull{2}}\quad\yiju{E2"第六到第八章只有少量填空/判断/选择，掌握基本概念＋公式套用即可"；E6 教材 8.8 节}\quad\nandu{中}

\begin{banfa}
\textbf{"看次数认应力"：}一次方（线性）$\to$黏性（分子尺度）；二次方$\to$湍流（微团尺度）。
题目给"$\tau=\mu\dd u/\dd y$"就是层流，给"$\rho l^2(\dd u/\dd y)^2$"就是湍流。
\textbf{分层回答"边界层内切应力如何分布"：}紧贴壁面的\textbf{黏性底层}里 $\tau_t$ 很小、黏性切应力主导；
越往外湍流切应力越大，在\textbf{湍流核心区}里湍流切应力远大于黏性切应力；
总切应力近似为常数（$\approx\tau_w$）。
\end{banfa}
\end{kaodian}

\section{湍流模型}

\begin{kaodian}{1}{湍流模型的分类与基本思想}
\textbf{是什么：}湍流模型的作用是\textbf{补充雷诺应力与平均速度场之间的关系，
使雷诺方程组封闭}。按"需要在方程组中补充几个新的微分方程"来分类，
可分为零方程、一方程、两方程、雷诺应力方程模型，以及不封闭而是直接求解的
大涡模拟（LES）与直接数值模拟（DNS）。

\begin{center}
\small
\begin{tabularx}{\textwidth}{@{}l>{\raggedright\arraybackslash}X>{\raggedright\arraybackslash}X>{\raggedright\arraybackslash}X@{}}
\toprule
类别 & 基本思想 & 优缺点 & 适用场合 \\
\midrule
零方程（混合长度） &
直接用代数式把 $\tau_t$ 与 $\dd\bar u/\dd y$ 挂钩，不引入新微分方程 &
优：最简单、快、易收敛；缺：忽略历史效应与逆压梯度，通用性最差 &
简单剪切流动、附着边界层、工程快速估算 \\
一方程（$k$ 方程） &
引入湍动能 $k$ 的输运方程求解 $k$，再配经验特征长度 $l$ &
优：考虑了湍能的历史/对流/扩散；缺：$l$ 仍需经验给定，通用性中等 &
中等复杂度的附着流、剪切流 \\
两方程（$k$-$\varepsilon$、$k$-$\omega$ SST） &
同时求解 $k$ 与湍流耗散率 $\varepsilon$（或比耗散率 $\omega$）两个输运方程，涡黏由二者组合得到 &
优：适用面最广、工程精度最好，是目前湍流工程计算的主流；
缺：计算量较大，对近壁处理与网格质量敏感 &
几乎全部工程湍流问题；\emph{SST} 版本对逆压梯度与分离流更好 \\
雷诺应力方程模型（RSM） &
直接求解 6 个雷诺应力分量的输运方程，不再用涡黏假设 &
优：能反映雷诺应力的各向异性、旋转与强曲率效应，精度高；
缺：方程多（含 6 个输运方程）、计算量大、不易收敛 &
强旋流、强曲率、分离与二次流 \\
大涡模拟（LES） &
用滤波把流动分为大尺度涡与小尺度涡：\emph{大涡直接求解}，小涡用\emph{亚格子模型}统一处理 &
优：能捕捉非定常大尺度结构，精度远超雷诺时均；
缺：网格与计算量巨大，近壁需特殊处理 &
非定常分离、尾涡、绕流噪声、燃烧等 \\
直接数值模拟（DNS） &
\emph{不加任何模型}，直接数值求解 N-S 方程，分辨\emph{全部}尺度的涡（网格 $\lesssim$ 最小涡尺度） &
优：理论上最精确，是研究湍流机理与检验模型的"标准答案"；
缺：计算量随 $\Rey$ 极高次幂增长，目前只能算低 $\Rey$、简单几何 &
基础湍流研究、简单几何低 $\Rey$ 流动 \\
\bottomrule
\end{tabularx}
\end{center}

\textbf{一句总括（很可能是简答题的答案骨架）：}
"湍流模型＝用\textbf{低阶的、可计算的量}（平均速度、$k$、$\varepsilon$、$\omega$）去\emph{近似表达}
\textbf{高阶的、未知的雷诺应力}，从而封闭时均方程组；
模型越复杂，能反映的物理机制越多、精度越高，但计算量与不确定参数也越大——
故工程上按'精度够用、成本可承受'来选型。"

\textbf{教材式号（8.9 节，书 p180--182）：}（教材式 8.78）雷诺应力的涡黏模型改写式
$\overline{u'_iu'_j}=\nu_t\left(\partial u_i/\partial x_j+\partial u_j/\partial x_i\right)-\dfrac23k\delta_{ij}$
（$\delta_{ij}$ 为克罗内克尔符号，$i=j$ 取 1、$i\neq j$ 取 0；$k$ 为湍动动能）；
（教材式 8.79）柯尔莫格罗夫-普朗特表达式 $\nu_t=C\sqrt{k}\,L$（$C$ 为经验常数、$L$ 为特征尺度）；
（教材式 8.80）$\nu_t=C_\mu k^2/\varepsilon$；（教材式 8.81）$k$-$\varepsilon$ 方程（含 $k$ 方程与 $\varepsilon$ 方程）。
\textbf{教材原话（8.9 节开头，书 p179，"不封闭"的标准论证）：}"湍流运动的控制方程包括 1 个时均连续性方程式（8.64）
和 3 个雷诺方程式（8.70），未知量包括 3 个时均速度分量、1 个时均压强和 6 个雷诺应力，共 10 个，
超过了方程的数目，因此方程组\textbf{不封闭}无法求解。"
\textbf{教材对一方程模型的评价（原话，书 p182）：}"由于 $k$ 方程模型考虑了湍动动能的迁移和扩散的传递以及湍流速度尺度的历史影响，
较混合长度模型优越。但应用此模型只限于简单的剪切层，因为复杂的流动，从经验去确定特征长度的分布很困难。"

\begin{tishi}
\textbf{订正（原表述 $\to$ 新表述）：}原稿把二方程模型写作"$k$-$\varepsilon$、$k$-$\omega$ SST"并列。
核对教材：\textbf{教材 8.9.4 只讲 $k$-$\varepsilon$ 二方程模型}，全章\textbf{未出现} $k$-$\omega$ 或 SST。
新表述：\emph{教材口径}为"二方程模型＝在 $k$ 方程之外再加上一个确定特征长度 $L$ 的偏微分方程，
应用最广泛的是 $k$-$\varepsilon$ 方程模型"（教材原话，书 p182）；本讲义保留 $k$-$\omega$ SST 作为\emph{工程扩展}，
但答题以教材口径为准。
另：教材 8.9 的四小节依次为"涡黏性模型（8.9.1）$\to$混合长度模型（8.9.2）$\to$一方程模型—$k$ 方程模型（8.9.3）$\to$二方程模型—$k$-$\varepsilon$ 方程模型（8.9.4）"，
教材\textbf{未涉及}雷诺应力方程模型（RSM）、大涡模拟（LES）、直接数值模拟（DNS）——本讲义这三项属\emph{扩展内容}，
\textbf{不标教材式号}，答题时若无把握可只说"教材介绍了零方程（混合长度）、一方程、二方程三类模型"。
\end{tishi}

\textbf{常考题型：}填空（"两方程模型中最常用的是\underline{\ $k$-$\varepsilon$ 模型\ }"；
"只求解大尺度涡、把小于网格尺度的涡用亚格子模型处理的方法称为\underline{\ 大涡模拟\ }"）；
选择（"不需要任何湍流模型的模拟方法是"$\to$DNS）；
简答（"简述湍流模型的分类及其基本思想"）。
\textbf{重要度 \xstarfull{1}}\quad\yiju{E2"第六到第八章只有少量填空/判断/选择"（最边缘内容）；E6 教材 8.9 节（简介）}\quad\nandu{中}

\begin{banfa}
\textbf{用"数方程个数"记分类：}\emph{不补}微分方程＝零方程（混合长度）；
补 1 个（$k$）＝一方程；补 2 个（$k$、$\varepsilon$ 或 $k$、$\omega$）＝两方程（工程主流）；
直接补 6 个＝雷诺应力方程模型。
\textbf{认名字：}$k$-$\varepsilon$、$k$-$\omega$ \textbf{SST} 常出现在选择题中，
"LES＝大涡模拟""DNS＝直接数值模拟"是必认的英文缩写。
\textbf{给分点：}答简答题时只要按"零方程→一方程→两方程→雷诺应力→LES→DNS"的顺序排开，
并各写一句"基本思想"，就已经覆盖了全部得分点。
\end{banfa}

\begin{yicuo}
\textbf{易错 7：}DNS 与 LES 的区别常被写反——\textbf{DNS 不加任何模型}（必须分辨所有尺度，
故网格要求最苛刻）；\textbf{LES 加亚格子模型}（只解析大涡，小涡被模化）。
记住"DNS＝全解、LES＝解大模小"。
另一个坑：把"$k$-$\varepsilon$ 模型"说成"两方程\emph{零}方程模型"（零方程模型是混合长度理论，没有微分方程）。
\end{yicuo}
\end{kaodian}

\section{本章小结}

\begin{center}
\begin{tabularx}{\textwidth}{@{}lccX@{}}
\toprule
考点 & 重要度 & 难度 & 一句话抓住 \\
\midrule
黏性流体的应力与本构关系 & \xstarfull{3} & 易 & 9 个分量、对称、独立 6 个；法向应力含 $-p$ 与 $2\mu\partial u_x/\partial x$ \\
N-S 方程与各项意义 & \xstarfull{3} & 中 & 惯性—压力—黏性—质量力；惯性项非线性故难解 \\
N-S 方程的解析解 & \xstarfull{2} & 易 & 库埃特流＝纯剪切流＋泊肃叶流；简单流动才有解析解 \\
边界层的定义与三重含义 & \xstarfull{3} & 易 & 薄、陡、旋；$\delta/l\sim1/\sqrt{\Rey_l}$ \\
边界层厚度 $\delta$ 的定义 & \xstarfull{3} & 易 & $u=0.99U_\infty$ 处；$\delta\propto\sqrt{\nu x/U}$ \\
两种边界层与转捩点 & \xstarfull{3} & 易 & 平板按 $\Rey_x$，临界 $5\times10^5$（管内按 $\Rey_d$，2000） \\
边界层内压强分布 & \xstarfull{3} & 中 & $\partial p/\partial y\approx0$，$p=p(x)$，由外部势流定 \\
边界层的适用条件 & \xstarfull{2} & 易 & 大 $\Rey$、未分离；尾流/分离/激波/低 $\Rey$ 不适用 \\
量级比较与普朗特方程 & \xstarfull{3} & 中 & 只留 $\nu\partial^2u/\partial y^2$；$\partial p/\partial y=0$；抛物型可推进求解 \\
边界条件 & \xstarfull{3} & 易 & 壁面 $u=v=0$；外缘 $u=U(x)$；$\dd p/\dd x=-\rho U\,\dd U/\dd x$ \\
布拉修斯相似性解 & \xstarfull{2} & 难 & $2f'''+ff''=0$，$f''(0)=0.332$，只用于平板 \\
平板层流公式（$\delta$、$\tau_w$、$C_f$、$C_D$） & \xstarfull{2} & 中 & $5.0/\sqrt{\Rey_x}$、$0.664/\sqrt{\Rey_x}$、$1.328/\sqrt{\Rey_l}$ \\
\bottomrule
\end{tabularx}
\end{center}

% 这张小结表原来是一整张 24 行的 tabularx。tabularx 不可跨页，整表高约 810pt
% 而 \textheight 约 700pt，实测报 "Overfull \vbox (88.06944pt too high)"，
% 即表格后段被挤出纸面（约 9mm 会被裁掉）。此处按行数对半拆成两张表、第二张重印表头，
% 未删任何一行、未缩小字号。
\begin{center}
\begin{tabularx}{\textwidth}{@{}lccX@{}}
\toprule
考点（续） & 重要度 & 难度 & 一句话抓住 \\
\midrule
卡门动量积分方程 & \xstarfull{2} & 难 & 平板：$\tau_w=\rho U^2\,\dd\delta_2/\dd x$ \\
$\delta_1$、$\delta_2$、$H$ 的定义与意义 & \xstarfull{3} & 中 & 1 管流量、2 管动量、$H$ 管分离（$H$ 大易分离） \\
假定剖面近似解法步骤 & \xstarfull{2} & 难 & 选剖面$\to$算厚度$\to$求 $\tau_w$$\to$代卡门$\to$积分求阻力 \\
湍流边界层公式 & \xstarfull{2} & 中 & 1/7 次方律；$0.37/\Rey_x^{1/5}$；$0.074/\Rey_l^{1/5}$ \\
混合边界层 & \xstarfull{2} & 中 & 全湍流值再减 $A/\Rey_l$；层流段越长阻力越小 \\
边界层分离的机制与判据 & \xstarfull{3} & 中 & 逆压梯度＋黏性；分离点 $\partial u/\partial y|_{y=0}=0$ \\
形状阻力与两种阻力对比 & \xstarfull{3} & 易 & 摩擦阻力 $\propto$湿面积；形状阻力 $\propto$迎流面积＋尾涡 \\
减阻措施 & \xstarfull{3} & 易 & 流线型、抽吸、提前转捩、微沟槽、高分子聚合物 \\
雷诺时均法与雷诺方程 & \xstarfull{3} & 中 & $\overline{u'}=0$、$\overline{u'^2}>0$；多出 $-\rho\overline{u'_iu'_j}$，方程组不封闭 \\
混合长度与湍流切应力 & \xstarfull{2} & 中 & $\tau_t=\rho l^2(\dd\bar u/\dd y)^2$；$l=\kappa y$ \\
湍流模型分类与思想 & \xstarfull{1} & 中 & 零$\to$一$\to$两方程$\to$雷诺应力$\to$LES$\to$DNS；DNS 不加模型 \\
\bottomrule
\end{tabularx}
\end{center}

\textbf{教材式号索引（本章出现的教材式号 $\to$ 教材节次；考前只需认"式号在哪节"）：}
\begin{center}
\begin{tabularx}{\textwidth}{@{}lXX@{}}
\toprule
教材式号 & 内容 & 教材节次 \\
\midrule
8.1$\sim$8.7 & 以应力表示的运动方程、牛顿内摩擦定律、广义牛顿内摩擦定律、平均压强与附加压应力 & 8.1.1$\sim$8.1.3 \\
8.8a$\sim$8.8c & 不可压缩均质黏性流体运动微分方程（N-S 方程） & 8.1.4 \\
8.9a、8.9b／8.10a、8.10b & 排挤厚度 $\delta^*$／动量损失厚度 $\theta$ & 8.2.2 \\
8.11 & 边界层雷诺数 $\Rey_x=U_0x/\nu$ & 8.2.3 \\
8.12$\sim$8.15 & 平面恒定边界层方程、量级分析、普朗特边界层方程及其另一种形式 & 8.3 \\
8.16 & $\delta\propto\sqrt{\nu x/U_0}$ & 8.3 \\
8.17$\sim$8.31 & 边界层控制方程与边界条件、相似性解、$f'''+\alpha ff''+\beta(1-f'^2)=0$、布拉修斯方程及各分量式 & 8.4 \\
8.32$\sim$8.35 & $C_f=0.664/\sqrt{\Rey_x}$、$D$、$C=1.32/\sqrt{\Rey_L}$、$\delta=5.0x/\sqrt{\Rey_x}$ & 8.4.3 \\
8.36$\sim$8.41 & 动量积分方程的推导与\textbf{卡门动量积分方程}（8.41） & 8.5 \\
8.42$\sim$8.50 & 平板边界层基本方程、层流"抛物线剖面"近似解（$5.477$、$1.46$） & 8.6.1 \\
8.51$\sim$8.56 & 1/7 次方律、$0.0225$、$\delta=0.37x/\Rey_x^{1/5}$、$0.0296$、$D$、$C=0.072/\Rey_L^{1/5}$ & 8.6.2 \\
8.57$\sim$8.59b & 混合边界层摩擦阻力及其摩阻系数 $0.074/\Rey_L^{1/5}-A/\Rey_L$ & 8.6.3 \\
8.60$\sim$8.62 & 摩擦阻力、压差阻力、绕流阻力 & 8.7.2 \\
8.63 & 圆球自由沉降速度 & 8.7.2 \\
8.64a、8.64b & 湍流时均连续性方程 & 8.8.2 \\
8.65$\sim$8.70 & 雷诺方程（8.70 为三个方向的雷诺方程） & 8.8.3 \\
8.71$\sim$8.73 & 湍动切应力与涡黏性模型 & 8.8.4、8.9.1 \\
8.74a$\sim$8.77 & 混合长度模型、$l=ky$、总切应力、对数分布 & 8.9.2 \\
8.78$\sim$8.81 & 涡黏模型改写、$\nu_t=C\sqrt{k}L$、$\nu_t=C_\mu k^2/\varepsilon$、$k$-$\varepsilon$ 方程 & 8.9.3、8.9.4 \\
\bottomrule
\end{tabularx}
\end{center}

\textbf{教材例题索引（本章 4 道，均为计算题同型题）：}例 8.1（书 p156--157）两无限长水平平板间恒定层流求速度分布 $\to$ 8.1 节；
例 8.2（书 p171）平板一面摩擦阻力与边界层厚度（用 $C_f=1.46/\sqrt{\Rey_L}$）$\to$ 8.6.1 节；
例 8.3（书 p171--172）平板两种放法摩擦阻力比 $F_1/F_2=0.908$ $\to$ 8.6.3 节（对应 E4 2015 年"平板摩擦阻力比"）；
例 8.4（书 p175）煤粉颗粒在烟气中能否沉降（$u_f=0.278$ m/s $<u=0.5$ m/s，被带走）$\to$ 8.7.2 节。
\textbf{本次增补的总原则（供复习时放心）：}全章以\textbf{教材（赵琴主编）为唯一准据}；
凡与教材不一致处已在原地用 \texttt{tishi} 写明"原表述 $\to$ 新表述"，凡判为教材自身印误处一律\textbf{两说并列}
（（教材式 8.34）的 $1.32$ 与 $1.328$、布拉修斯博士论文的 $1980$ 年与 $1908$ 年），\emph{不擅自删原说法}。
